\BourbakiAppendixExercises{习题}

\begin{enumerate}
\def\labelenumi{\arabic{enumi})}
\tightlist
\item
  设 \(\mathrm{A}\) 是一个 \(k\)-代数。
\end{enumerate}

\begin{enumerate}
\def\labelenumi{\alph{enumi})}
\item
  证明：包含 AA 的正则（左或右）理想只有 A 本身。
\item
  设 \(\mathfrak{a}\) 是 \(A\) 的一个双边理想。它既作为左理想又作为右理想都是正则的，当且仅当 \(k\)-代数 \(A / \mathfrak{a}\) 有单位元。
\end{enumerate}

\begin{enumerate}
\def\labelenumi{\arabic{enumi})}
\setcounter{enumi}{1}
\tightlist
\item
  设 \(\mathbf{A}\) 是一个 \(k\)-代数，\(\mathfrak{a}\) 与 \(\mathfrak{b}\) 是 \(\mathbf{A}\) 的正则（左）理想。
\end{enumerate}

\begin{enumerate}
\def\labelenumi{\alph{enumi})}
\item
  理想 \(a \cap b\) 是正则的，当且仅当存在模 a 的右单位元 u 与模 b 的右单位元 v，使得 \(u - v \in a + b\)。
\item
  由此推出：若 a 是双边理想且是右正则的，或者若 a 是极大的，则 \(a \cap b\) 是正则的。c) 设 K 是一个交换域，L 是两个变元 X 与 Y 上的自由结合 K-代数（III, §2, No.~7, 第 448 页），而 A 是 L 中由次数 \(\geqslant 1\) 的元素组成的 K-子代数。证明 A 的理想 \(A(1 - X)\) 与 \(A(1 - Y)\) 都是正则的，但它们的交约化为 0。
\end{enumerate}

\begin{enumerate}
\def\labelenumi{\arabic{enumi})}
\setcounter{enumi}{2}
\tightlist
\item
  设 \(\mathbf{A}\) 是一个 \(k\)-代数，\(\mathfrak{a}\) 是 \(\mathbf{A}\) 的一个正则左理想，\(u\) 与 \(v\) 是模 \(\mathfrak{a}\) 的右单位元。\(a)\) 给出一个 \(u - v\) 不属于 \(\mathfrak{a}\) 的例子（参见 III, §9, 第 644 页，习题 1）。由此推出一个使映射 \(\widetilde{\mathfrak{a}} \mapsto \widetilde{\mathfrak{a}} \cap \mathbf{A}\)（VIII, 第 437 页，命题 3）不是单射的例子。
\end{enumerate}

\begin{enumerate}
\def\labelenumi{\alph{enumi})}
\setcounter{enumi}{1}
\tightlist
\item
  假设理想 \(\mathfrak{a}\) 是双边理想，且 \(k\)-代数 \(B = A / \mathfrak{a}\) 不包含任何非零双边理想 \(\mathfrak{b}\) 使得 \(B\mathfrak{b} = 0\)。证明 \(u - v\) 属于 \(\mathfrak{a}\)。特别地，当 \(A\) 交换时就是这种情形。
\end{enumerate}

\begin{enumerate}
\def\labelenumi{\arabic{enumi})}
\setcounter{enumi}{3}
\tightlist
\item
  设 A 是一个 \(k\)-代数，\(\mathfrak{a}\) 是 A 的一个左理想；我们把 \(\mathfrak{a}\) 视为一个 \(k\)-代数。证明我们有 \(\Re(A) \cap \mathfrak{a} \subset \Re(\mathfrak{a})\)（应用 VIII, 第 441 页，命题 7），并且当 \(\mathfrak{a}\) 是双边理想时有等式（注意 A 上的每个单伪模在 \(\mathfrak{a}\) 上是单的或被 \(\mathfrak{a}\) 零化）。给出一个左理想 \(\mathfrak{a}\) 的例子，使得 \(\Re(A) \cap \mathfrak{a} \neq \Re(\mathfrak{a})\)（取 A 为一个矩阵代数）。
\end{enumerate}

¶ 5) 假设 A 上的每个伪模 M 都是 AM 与 M 中由被 A 零化的元素组成的子伪模的直和。证明 A 有单位元。（取 \(M = \widetilde{A}_{s}\)，证明 A 容许一个右单位元 u。取 \(M = A_{s}\)，证明对 A 的每个非零元素 x 我们有 \(Ax \neq 0\)。由此断定 u 也是一个左单位元。）

\begin{enumerate}
\def\labelenumi{\arabic{enumi})}
\setcounter{enumi}{5}
\tightlist
\item
  我们称一个 k-代数是本原的，如果它容许一个既单又忠实的伪模（即其零化子约化为 0 的伪模）。
\end{enumerate}

\begin{enumerate}
\def\labelenumi{\alph{enumi})}
\item
  证明：若代数 \(\tilde{\mathbf{A}}\) 是本原的，则 \(\mathbf{A}\) 也是本原的。
\item
  若代数 A 是本原的且有单位元，则代数 \(\widetilde{A}\) 不是本原的。
\item
  假设 k 是一个体，且代数 A 是本原的但没有单位元。证明 \(\widetilde{A}\) 是本原的（证明：若存在 A 上的一个忠实伪模，它作为 \(\widetilde{A}\)-模不忠实，则 A 容许一个右单位元 u，并由此推出矛盾）。
\end{enumerate}

\begin{enumerate}
\def\labelenumi{\arabic{enumi})}
\setcounter{enumi}{6}
\item
  我们称 \(k\)-代数 A 是阿廷的，如果 A 的左理想的每个递减序列都是平稳的。环 \(\widetilde{\mathrm{A}}\) 是阿廷的，当且仅当代数 A 与环 \(k\) 都是阿廷的。证明：若 \(k\) 是一个体，则无根的阿廷代数 A 有单位元（注意代数 \(\widetilde{\mathrm{A}}\) 是半单的）。
\item
  假设 \(k\) 是一个体，且 \(A\) 是 \(k\) 上的一个本原代数。证明：若 \(A\) 容许一个在 \(k\) 上有限维的非零左理想，则代数 \(A\) 在 \(k\) 上是有限维的，从而是一个单环（参见习题 7）。
\item
  假设 k 是一个体，且 A 是一个幂零 k-代数，即存在一个整数 n 使得 \(A^{n}=0\)。证明 A 在 k 上是有限维的（对使 \(A^{n}=0\) 的最小整数 n 用归纳法）。
\item
  假设 \(k\) 是一个体，\(A\) 是一个无限维 \(k\)-代数，并且 \(A\) 的每个真左理想都在 \(k\) 上有限维。
\end{enumerate}

\begin{enumerate}
\def\labelenumi{\alph{enumi})}
\item
  证明我们有 \(Aa = 0\) 对每个左理想 \(a \neq A\) 成立（考虑同态 \(\gamma : A \to \operatorname{End}_{k}(a)\)）。
\item
  设 \(\mathfrak{r}\) 是 A 的极大左理想之和。证明 \(\mathfrak{r}\) 异于 A（利用习题 9），\(\mathfrak{r}\) 是一个双边理想，且 \(A / \mathfrak{r}\) 是一个体（参见习题 7）。c) 设 \(u\) 是模 \(\mathfrak{r}\) 的一个（左或右）单位元。证明我们有 \(\mathfrak{ru} = \mathfrak{r}\)（注意我们有 \(Au = A\)）。
\item
  证明 rA 为零（考虑 r 在 A 中的右零化子，并按 a) 的方式推理），然后证明 r 为零。由此得出 A 是一个体。
\end{enumerate}

¶ 11) 假设 k 是一个体，A 是一个 k-代数，并且 A 的子代数的每个递减序列都是平稳的。

\begin{enumerate}
\def\labelenumi{\alph{enumi})}
\item
  证明代数 \(\mathbf{A}\) 是阿廷的，且其根在 \(k\) 上有限维（应用习题 9）。
\item
  设 S 是 A/ \(\Re(A)\) 的一个单分量（习题 7）；它同构于一个体 D 上的矩阵代数，其中心为 Z。证明 D 在 Z 上有限维（若 D 含有一个在 Z 上超越的元素，构造 D 中一个无穷严格递减的子体序列；若 D 在 Z 上是代数的且无限维，构造 D 中一个无穷严格递增的子体序列，并考虑它们在 D 中的交换子）。
\item
  证明：若 Z 在 k 上是无限维的，则我们有 S = D（考虑 B 的幂零子代数）。
\item
  证明 Z 是 k 的一个代数扩张。设 p 是 k 的特征指数，L 是 Z 所含的 k 的最大可分扩张。证明 \(L \cap Z^{p}\) 在 \(L^{p}\) 上有有限次数（考虑 \(L \cap Z^{p}\) 在 \(L^{p}\) 上的一个 p-基）。
\item
  L 的子扩张的每个递减序列都是平稳的。给出一个例子说明这个条件并不蕴含 L 在 k 上有有限次数（参见 V, §12, 第 167 页，习题 3）。
\item
  反之，设 B 是一个满足 a) 的条件的 k-代数，使得 \(\mathrm{B}/\Re(\mathrm{B})\) 的每个单分量都满足 b)、c)、d) 与 e) 的条件。证明 B 的子代数的每个递减序列都是平稳的。
\end{enumerate}

（归结为 B 是一个体且中心为 Z 的情形。设 \((\mathrm{D}_{n})\) 是 B 的子体的一个递减序列；我们可以假设由 \(D_{n}\) 与 Z 生成的体 D 与 n 无关。注意我们可以写 \(D = E \otimes_{L} Z\)，其中 L 是 Z 的一个在 k 上有有限次数的子体，而 E 是一个在 L 上有有限次数的体，其中心包含 L。证明 \(D_{n}\) 这时就是由 E 与 \(D_{n} \cap Z\) 生成的体。于是我们归结为 B 是一个交换域且是 k 的 p-根式扩张的情形。若 \((\mathrm{F}_{n})\) 是 B 的子扩张的一个递减序列，考虑子体 \(k(\mathrm{F}_{n}^{p^{\infty}})\)，其中 \(F_{n}^{p^{\infty}} = \cap_{f} F_{n}^{p^{f}}\)。）

¶ 12) a) 假设 k 是一个体，k-代数 A 是阿廷的，并且 A 的子代数的每个递增序列都是平稳的。证明 A 在 k 上有有限次数。（按习题 11 的方式推理。注意体 \(k(\mathrm{X})\) 含有一个无穷严格递增的子代数序列，这利用了 VII, §1, No.~5, 第 6 页，命题 6。）b) 证明在多项式代数 \(k[X]\) 中，子代数的每个递增序列都是平稳的。（设 \(f \in k[X]\) 是一个非常数多项式，A 是由 1 与 f 生成的 \(k[X]\) 的 k-子代数。注意 A 是一个主理想整环，且 \(k[X]\) 是一个有限生成 A-模。）

\begin{enumerate}
\def\labelenumi{\arabic{enumi})}
\setcounter{enumi}{12}
\item
  假设 \(k\) 是一个体，且 \(A\) 容许一个由幂零元素组成的在 \(k\) 上的有限基。证明 \(A\) 是幂零的（习题 9；归结为代数 \(A\) 是半单的情形（参见习题 7），然后通过扩张基域归结为形如 \(\mathbf{M}_n(\mathrm{K})\) 的情形）。
\item
  设 A 与 B 是 k-代数；(A,B)-伪双模是一个集合 E，它赋有 A 上的左伪模结构与 B 上的右伪模结构，且使得底 k-模结构一致。这对应于赋予 E 一个 \(\widetilde{A}\) 与 \(\widetilde{B}\) 这两个代数上的双模结构。E 的一个乘子是一个偶对 (S,D)，其中 S:B→E 是右 B-伪模的同态，D:A→E 是左 A-伪模的同态，满足 \(a\mathrm{S}(b)=\mathrm{D}(a)b\)（对 \(a\in A\)、\(b\in B\)）。用 \(\mathcal{M}(\mathrm{E})\) 表示 E 的乘子所成的集合；它容许一个自然的 k-模结构。
\end{enumerate}

\(a\) ) 对每个 \(x\in \mathrm{E}\)，用 \(\mathrm{D}_x\)（相应地，\(\mathrm{S}_x\)）表示映射 \(a\mapsto ax\)（相应地，\(b\mapsto xb\)）。证明偶对 \(m_{x} = (\mathrm{D}_{x},\mathrm{S}_{x})\) 是 E 的一个乘子。

\begin{enumerate}
\def\labelenumi{\alph{enumi})}
\setcounter{enumi}{1}
\tightlist
\item
  对每个 \(m = (\mathrm{S}, \mathrm{D}) \in \mathcal{M}(\mathrm{E})\)、\(a \in A\) 与 \(b \in B\)，置 \(am = m_{\mathrm{D}(a)}\) 与 \(mb = m_{\mathrm{S}(b)}\)。证明这在 k-模 \(\mathcal{M}(\mathrm{E})\) 上定义了一个 (A, B)-伪双模结构，且映射 \(x \mapsto m_x\) 是 (A, B)-伪双模的一个同态。
\end{enumerate}

\begin{enumerate}
\def\labelenumi{\arabic{enumi})}
\setcounter{enumi}{14}
\tightlist
\item
  设 \(\mathrm{A}\) 是一个 \(k\)-代数。\(\mathrm{A}\) 的乘子是 (A,A)-伪双模 \(\mathrm{A}\) 的一个乘子。乘子所成的 \(k\)-代数（\(\mathrm{A}\) 的）用 \(\mathcal{M}(\mathrm{A})\) 表示。
\end{enumerate}

\begin{enumerate}
\def\labelenumi{\alph{enumi})}
\item
  证明映射 \(\mu: a \mapsto m_{a}\) 是一个 k-代数同态，且我们有 \(am = m_{a}m\) 与 \(ma = mm_{a}\)（对 \(a \in A, m \in \mathcal{M}(A)\)）。因此同态 \(\mu\) 的像是环 \(\mathcal{M}(A)\) 的一个双边理想。
\item
  k-代数 A 的一个双边理想 b 称为忠实的，如果它的左零化子与右零化子都为零。证明：若 b 在 A 中是忠实的，则同态 \(\mu\) 是单射，且理想 \(\mu(\mathrm{A})\) 在 \(\mathcal{M}(\mathrm{A})\) 中是忠实的。
\item
  设 B 是一个包含 A 作为双边理想的有单位元的 k-代数。对 \(b \in B\)，设 \(\sigma_{b}\) 与 \(\delta_{b}\) 是从 A 到 A 的映射，使得对 \(a \in A\) 有 \(\sigma_{b}(a) = ba\) 与 \(\delta_{b}(a) = ab\)，并置 \(\pi(b) = (\sigma_{b}, \delta_{b})\)。证明 \(\pi(b) \in \mathcal{M}(A)\) 对每个 \(b \in B\) 成立，且如此得到的映射 \(\pi : B \to \mathcal{M}(A)\) 是一个环同态，它在 A 上的限制是映射 \(a \mapsto m_{a}\)。证明：若理想 A 在 B 中是忠实的，则 \(\pi\) 是单射。
\end{enumerate}

\chapter{非交换体上的行列式}

在本附录中，D 是一个体，而 \(D_{ab}^{*}\) 是 D 的乘法群 \(D^{*}\) 对其导群的商（I, §6, No.~2, 第 70 页）。群 \(D_{ab}^{*}\) 是交换的，且当体 D 交换时，\(D_{ab}^{*}\) 可与 \(D^{*}\) 等同。从 \(D^{*}\) 到 \(D_{ab}^{*}\) 的典范同态用 \(\pi\) 表示。

\section{交错多重线性形式的一个推广}

设 \(\mathrm{V}\) 是体 \(\mathrm{D}\) 上的一个右向量空间，维数有限，为 \(n \geqslant 0\)。用 \(\mathrm{B(V)}\) 表示 \(\mathrm{V}\) 的基所成的集合，用 \(\Omega(\mathrm{V})\) 表示满足下列两个条件的映射 \(\omega\)（从 \(\mathrm{B(V)}\) 到 \(\mathrm{D}_{\mathrm{ab}}^{*}\)）所成的集合：

\begin{enumerate}
\def\labelenumi{\alph{enumi})}
\tightlist
\item
  若 \(\lambda_1, \ldots, \lambda_n\) 是 \(D^*\) 中的元素，则我们有
\end{enumerate}

\[
\omega (v _ {1} \lambda_ {1}, \ldots , v _ {n} \lambda_ {n}) = \pi (\lambda_ {1} \dots \lambda_ {n}) \omega (v _ {1}, \ldots , v _ {n})\tag{1}
\]

对 V 的每个基 \((v_{1},\ldots ,v_{n})\) 成立。

\begin{enumerate}
\def\labelenumi{\alph{enumi})}
\setcounter{enumi}{1}
\tightlist
\item
  若 \(i\) 与 \(j\) 是区间 \([1, n]\) 中两个不同的整数，则我们有
\end{enumerate}

\[
\omega (v _ {1}, \ldots , v _ {i - 1}, v _ {i} + v _ {j}, v _ {i + 1}, \ldots , v _ {n}) = \omega (v _ {1}, \ldots , v _ {n})\tag{2}
\]

对 V 的每个基 \((v_{1},\ldots ,v_{n})\) 成立。

设 \(\omega\) 是 \(\Omega(\mathrm{V})\) 的一个元素。设 \((v_1, \ldots, v_n)\) 是 \(\mathrm{V}\) 的一个基，\(i\) 是区间 \([1, n-1]\) 中的一个整数；由性质 b)，我们有

\[
\begin{array}{r l} & {\omega (v _ {1}, \ldots , v _ {i}, v _ {i + 1}, \ldots , v _ {n}) = \omega (v _ {1}, \ldots , v _ {i} + v _ {i + 1}, v _ {i + 1}, \ldots , v _ {n})} \\ & {\qquad = \omega (v _ {1}, \ldots , v _ {i} + v _ {i + 1}, v _ {i + 1} - (v _ {i} + v _ {i + 1}), \ldots , v _ {n})} \\ & {\qquad = \omega (v _ {1}, \ldots , v _ {i} + v _ {i + 1}, - v _ {i}, \ldots , v _ {n})} \\ & {\qquad = \omega (v _ {1}, \ldots , (v _ {i} + v _ {i + 1}) - v _ {i}, - v _ {i}, \ldots , v _ {n}),} \end{array}
\]

从而

\[
\omega (v _ {1}, \dots , v _ {i}, v _ {i + 1}, \dots , v _ {n}) = \pi (- 1) \omega (v _ {1}, \dots , v _ {i + 1}, v _ {i}, \dots , v _ {n}).\tag{3}
\]

由于对称群 \(\mathfrak{S}_n\) 由区间 \([1,n]\) 中两个相邻元素的对换生成（I, §5, No.~7, 第 63 页，命题 9），公式 (3) 可推广为

\[
\omega (v _ {\sigma (1)}, \ldots , v _ {\sigma (n)}) = \pi (\varepsilon (\sigma)) \omega (v _ {1}, \ldots , v _ {n}),\tag{4}
\]

其中 \(\sigma\) 属于 \(\mathfrak{S}_n\)，而 \(\varepsilon (\sigma)\) 是它的符号差（I, §5, No.~7, 第 64 页）。

公式 (2) 可如下推广。首先，对每个 \(\lambda\)（在 \(\mathrm{D}^*\) 中），我们有

\[
\begin{array}{r l} & {\omega (v _ {1}, \ldots , v _ {n}) = \pi (\lambda) \omega (v _ {1}, \ldots , v _ {i} \lambda^ {- 1}, \ldots , v _ {n})} \\ & {\qquad = \pi (\lambda) \omega (v _ {1}, \ldots , v _ {i} \lambda^ {- 1} + v _ {j}, \ldots , v _ {n})} \\ & {\qquad = \omega (v _ {1}, \ldots , (v _ {i} \lambda^ {- 1} + v _ {j}) \lambda , \ldots , v _ {n}),} \end{array}
\]

也就是说，

\[
\omega (v _ {1}, \dots , v _ {n}) = \omega (v _ {1}, \dots , v _ {i} + v _ {j} \lambda , \dots , v _ {n}).\tag{5}
\]

由一个直接的归纳法，我们推出

\[
\omega (v _ {1}, \ldots , v _ {n}) = \omega (v _ {1}, \ldots , v _ {i} + w, \ldots , v _ {n})\tag{6}
\]

对每个线性组合 \(w = \sum_{j \neq i} v_j \lambda_j\)（诸向量 \(v_j\)，其中 j 在区间 \([1, n]\) 中且异于 i）成立。

\section{一个唯一性定理}

在本小节中，对任何严格正的整数 m、区间 \([1,m]\) 中的任何整数 i 以及任何序列 \((v_{1},\ldots,v_{m})\)（在 \(V^{m}\) 中），我们用 \((v_{1},\ldots,\widehat{v_{i}},\ldots,v_{m})\) 表示序列 \((v_{1},\ldots,v_{i-1},v_{i+1},\ldots,v_{m})\)（在 \(V^{m-1}\) 中）。

命题 1. --- 设 W 是 V 中的一个超平面，e 是 V---W 中的一个向量。设 \(\varphi\) 是 \(\Omega(\mathrm{W})\) 的一个元素。存在唯一的元素 \(\omega\) 属于 \(\Omega(\mathrm{V})\)，使得我们有

\[
\omega (w _ {1}, \ldots , w _ {n - 1}, e) = \varphi (w _ {1}, \ldots , w _ {n - 1})\tag{7}
\]

对 W 的每个基 \((w_{1},\ldots ,w_{n - 1})\) 成立。

\begin{enumerate}
\def\labelenumi{\Alph{enumi})}
\tightlist
\item
  \(\omega\) 的唯一性：设 \(\omega\) 是 \(\Omega(\mathrm{V})\) 的一个元素。设 \((v_{1},\ldots,v_{n})\) 是 V 的一个基，\(\mu_{1},\ldots,\mu_{n}\) 是 e 在这个基中的坐标。用 I 表示从 1 到 n（含端点）中使 \(\mu_{i}\neq0\) 的整数所成的集合；它非空。设 i 是 I 的一个元素。序列 \((v_{1},\ldots,\widehat{v_{i}},\ldots,v_{n},e)\) 是 V 的一个基，且由于 \(e-v_{i}\mu_{i}\) 是序列 \((v_{1},\ldots,\widehat{v_{i}},\ldots,v_{n})\) 的一个线性组合，公式 (6) 蕴含
\end{enumerate}

\[
\omega (v _ {1}, \ldots , \widehat {v _ {i}}, \ldots , v _ {n}, v _ {i} \mu_ {i}) = \omega (v _ {1}, \ldots , \widehat {v _ {i}}, \ldots , v _ {n}, e).\tag{8}
\]

由公式 (1) 和 (3) 可知，这个等式的左边等于 \(\pi(-1)^{n-i}\pi(\mu_i)\omega(v_1,\ldots,v_n)\)。用 p 表示 V 的以 W 为像、以 eD 为核的投影。向量 \(v_j - p(v_j)\) 都与 e 成比例，反复应用公式 (6) 表明公式 (8) 的右边等于 \(\omega(p(v_1),\ldots,\widehat{p(v_i)},\ldots,p(v_n),e)\)。

由此得出，若 \(\omega\) 满足关系式 (7)，则我们有

\[
\omega (v _ {1}, \dots , v _ {n}) = \pi (- 1) ^ {n - i} \pi (\mu_ {i}) ^ {- 1} \varphi (p (v _ {1}), \dots , \widehat {p (v _ {i})}, \dots , p (v _ {n}))\tag{9}
\]

对 I 中的每个 i 成立，这就证明了 \(\omega\) 的唯一性。

\begin{enumerate}
\def\labelenumi{\Alph{enumi})}
\setcounter{enumi}{1}
\tightlist
\item
  \(\omega\) 的构造：设 \((v_{1},\ldots,v_{n})\) 是 V 的一个基；我们像上面那样定义 \(\mu_{1},\ldots,\mu_{n}\) 与 I。对 I 中的每个 \(i\)，假设 \(\mu_{i}\neq0\) 蕴含序列 \((v_{1},\ldots,\widehat{v_{i}},\ldots,v_{n},e)\) 生成空间 V，因此序列 \((p(v_{1}),\ldots,\widehat{p(v_{i})},\ldots,p(v_{n}))\) 生成空间 W = p(V)。换言之，后一序列是 W 的一个基，且
\end{enumerate}

\[
t _ {i} = \pi (- 1) ^ {n - i} \pi (\mu_ {i}) ^ {- 1} \varphi (p (v _ {1}), \ldots , \widehat {p (v _ {i})}, \ldots , p (v _ {n}))\tag{10}
\]

定义了 \(\mathrm{D}_{\mathrm{ab}}^{*}\) 的一个元素。

设 i 与 j 是 I 的两个满足 i \textless{} j 的元素；让我们证明等式 \(t_{i} = t_{j}\)。按定义，向量 \(v_{i}\mu_{i} + v_{j}\mu_{j} - e\) 是诸向量 \(p(v_{k})\)（k 异于 i 与 j）的一个线性组合。因此，\(p(v_{i})\mu_{i} + p(v_{j})\mu_{j}\) 是诸向量 \(p(v_{k})\)（k 异于 i 与 j）的一个线性组合。于是由公式 (6)，我们有

\[
\begin{array}{c} \varphi (p (v _ {1}), \ldots , \widehat {p (v _ {i})}, \ldots , \widehat {p (v _ {j})}, \ldots , p (v _ {n}), p (v _ {i}) \mu_ {i}) \\ = \varphi (p (v _ {1}), \ldots , \widehat {p (v _ {i})}, \ldots , \widehat {p (v _ {j})}, \ldots , p (v _ {n}), - p (v _ {j}) \mu_ {j}). \end{array}
\]

应用公式 (1) 与 (3)，我们推出

\[
\begin{array}{c} \varphi (p (v _ {1}), \ldots , \widehat {p (v _ {j})}, \ldots , p (v _ {n})) \pi (\mu_ {i}) \pi (- 1) ^ {n - i - 1} \\ = \varphi (p (v _ {1}), \ldots , \widehat {p (v _ {i})}, \ldots , p (v _ {n})) \pi (- \mu_ {j}) \pi (- 1) ^ {n - j}, \end{array}
\]

从而 \(t_i = t_j\)。

证明了这一点之后，我们把 \(\omega(v_1, \ldots, v_n)\) 定义为诸 \(t_i\)（\(i\) 属于 I）的公共值。这样我们就构造了一个映射 \(\omega\)（从 B(V) 到 D \(_{\text{ab}}^*\)），它满足关系式 (9)。设 \((w_1, \ldots, w_{n-1})\) 是 W 的一个基；若置 \(v_i = w_i\)（对 \(1 \leqslant i \leqslant n-1\)）且置 \(v_n = e\)，则我们有 I = \{n\} 与 \(\mu_n = 1\)。我们还有 \(p(v_i) = w_i\)（对 \(1 \leqslant i \leqslant n-1\)），而公式 \(\omega(w_1, \ldots, w_{n-1}, e) = \varphi(w_1, \ldots, w_{n-1})\) 是 (9) 的一个特殊情形。

现在我们必须证明 \(\omega\) 属于 \(\Omega (\mathrm{V})\)。

\begin{enumerate}
\def\labelenumi{\Alph{enumi})}
\setcounter{enumi}{2}
\tightlist
\item
  公式 (1) 的证明：设 \(\lambda_1, \ldots, \lambda_n\) 是 \(D^*\) 的元素，\((v_1, \ldots, v_n)\) 是 \(V\) 的一个基。我们像上面那样定义 \(\mu_1, \ldots, \mu_n\) 与 \(I\)。取定一个 \(i\) 属于 \(I\)。由于我们有 \(e = \sum_{j=1}^{n} (v_j \lambda_j)(\lambda_j^{-1} \mu_j)\)，公式 (9) 蕴含
\end{enumerate}

\[
\begin{array}{l l} (1 1) & \omega (v _ {1} \lambda_ {1}, \ldots , v _ {n} \lambda_ {n}) \\ & \qquad = \pi (- 1) ^ {n - i} \pi (\lambda_ {i} ^ {- 1} \mu_ {i}) ^ {- 1} \varphi (p (v _ {1}) \lambda_ {1}, \ldots , \widehat {p (v _ {i}) \lambda_ {i}}, \ldots , p (v _ {n}) \lambda_ {n}). \end{array}
\]

但我们有 \(\pi (\lambda_i^{-1}\mu_i)^{-1} = \pi (\mu_i)^{-1}\pi (\lambda_i)\)，且

\[
\begin{array}{c} \varphi (p (v _ {1}) \lambda_ {1}, \ldots , \widehat {p (v _ {i}) \lambda_ {i}}, \ldots , p (v _ {n}) \lambda_ {n}) \\ = \varphi (p (v _ {1}), \ldots , \widehat {p (v _ {i})}, \ldots , p (v _ {n})) \pi (\lambda_ {1} \ldots \widehat {\lambda_ {i}} \ldots \lambda_ {n}). \end{array}
\]

比较公式 (9) 与 (11) 就得到所要的关系式

\[
\omega (v _ {1} \lambda_ {1}, \ldots , v _ {n} \lambda_ {n}) = \omega (v _ {1}, \ldots , v _ {n}) \pi (\lambda_ {1} \ldots \lambda_ {n}).
\]

\begin{enumerate}
\def\labelenumi{\Alph{enumi})}
\setcounter{enumi}{3}
\tightlist
\item
  公式 (2) 的证明：设 \((v_{1},\ldots,v_{n})\) 是 V 的一个基，i、j 是区间 {[}1,n{]} 中两个不同的整数；我们像先前那样定义 \(\mu_{1},\ldots,\mu_{n}\)。考虑 V 的基 \((v_{1}^{\prime},\ldots,v_{n}^{\prime})\)，它由 \(v_{i}^{\prime}=v_{i}+v_{j}\) 与 \(v_{k}^{\prime}=v_{k}\)（\(k\neq i\)）定义，并引入 e 在这个基中的坐标 \(\mu_{1}^{\prime},\ldots,\mu_{n}^{\prime}\)；它们满足 \(\mu_{j}^{\prime}=\mu_{j}-\mu_{i}\) 与 \(\mu_{k}^{\prime}=\mu_{k}\)（\(k\neq j\)）。我们置 \(t=\omega(v_{1},\ldots,v_{n})\) 与 \(t^{\prime}=\omega(v_{1}^{\prime},\ldots,v_{n}^{\prime})\)。我们必须证明 t 与 \(t^{\prime}\) 相等。
\end{enumerate}

首先注意，按 \(\omega\) 的定义，我们有

\begin{enumerate}
\def\labelenumi{(\arabic{enumi})}
\setcounter{enumi}{11}
\tightlist
\item
\end{enumerate}

\[
t = \pi (- 1) ^ {n - k} \pi (\mu_ {k}) ^ {- 1} \varphi (p (v _ {1}), \ldots , \widehat {p (v _ {k})}, \ldots , p (v _ {n})),\tag{13}
\]

\[
t ^ {\prime} = \pi (- 1) ^ {n - k} \pi (\mu_ {k} ^ {\prime}) ^ {- 1} \varphi (p (v _ {1} ^ {\prime}), \ldots , \widehat {p (v _ {k} ^ {\prime})}, \ldots , p (v _ {n} ^ {\prime}))
\]

对每个 \(k\)（使 \(\mu_{k}\) 与 \(\mu_k^\prime\) 都非零）成立。

\begin{enumerate}
\def\labelenumi{\alph{enumi})}
\item
  若 \(\mu_i \neq 0\)，则序列 \((p(v_1), \ldots, \widehat{p(v_k)}, \ldots, p(v_n))\) 与 \((p(v_1'), \ldots, \widehat{p(v_k'}, \ldots, p(v_n'))\) 相等，且我们有 \(\mu_i = \mu_i'\)，从而 \(t = t'\)。
\item
  若存在一个指标 \(k\)（异于 \(i\) 与 \(j\)）使 \(\mu_k \neq 0\)，则我们有 \(p(v_l') = p(v_l)\)（对 \(l \neq i\)），且 \(p(v_i') = p(v_i) + p(v_j)\)。元素 \(p(v_i)\)
\end{enumerate}

与 \(p(v_{j})\) 都属于序列 \((p(v_{1}),\ldots ,\widehat{p(v_{k})},\ldots ,p(v_{n}))\)。由于 \(\varphi\) 属于 \(\Omega (\mathrm{W})\)，VIII, 第 447 页的公式 (2) 适用；它给出

\[
\varphi (p (v _ {1}), \dots , \widehat {p (v _ {k})}, \dots , p (v _ {n})) = \varphi (p (v _ {1} ^ {\prime}), \dots , \widehat {p (v _ {k} ^ {\prime})}, \dots , p (v _ {n} ^ {\prime})).
\]

但我们还有 \(\mu_{k} = \mu_{k}^{\prime}\)，从而 \(t = t'\)。

\begin{enumerate}
\def\labelenumi{\alph{enumi})}
\setcounter{enumi}{2}
\tightlist
\item
  只剩下考察指标 \(k\)（使 \(\mu_k \neq 0\) 的唯一指标）是 \(j\) 的情形。这时我们有 \(\mu_j' = \mu_j\)、\(e = v_j\mu_j\) 与 \(p(v_j) = 0\)。在公式 (12) 与 (13) 中取 \(k = j\)。由于序列 \((p(v_1), \ldots, \widehat{p(v_j)}, \ldots, p(v_n))\) 与 \((p(v_1'), \ldots, \widehat{p(v_j'), \ldots, p(v_n'))}\) 相等，我们有 \(t = t'\)。
\end{enumerate}

推论 1. --- 设 \((e_1, \ldots, e_n)\) 是 V 的一个基，t 是 \(\mathrm{D}_{\mathrm{ab}}^*\) 的一个元素。存在唯一的元素 \(\omega\) 属于 \(\Omega(\mathrm{V})\)，使得 \(\omega(e_1, \ldots, e_n) = t\)。

我们对 n 用归纳法证明该推论。若 n = 0，则 V 的唯一的基是空基，故断言为真。现在设 \(n \geqslant 1\)；用 W 表示由 \(e_{1}, \ldots, e_{n-1}\) 生成的 V 的子空间。由命题 1，存在一个双射 \(\Lambda\)（从 \(\Omega(\mathrm{V})\) 到 \(\Omega(\mathrm{W})\)），满足

\[
(\Lambda (\omega)) (w _ {1}, \ldots , w _ {n - 1}) = \omega (w _ {1}, \ldots , w _ {n - 1}, e _ {n})
\]

对 W 的每个基 \((w_{1},\ldots,w_{n-1})\) 与每个 \(\omega\)（属于 \(\Omega(\mathrm{V})\)）成立。由归纳假设，存在唯一的元素 \(\varphi\) 属于 \(\Omega(\mathrm{W})\)，使得 \(\varphi(e_{1},\ldots,e_{n-1})=t\)。关系式 \(\omega(e_{1},\ldots,e_{n})=t\)（\(\omega\) 属于 \(\Omega(\mathrm{V})\)）等价于 \(\Lambda(\omega)=\varphi\)。由此即得推论 1。

推论 2. --- 设 \(\omega\) 与 \(\omega'\) 是 \(\Omega(V)\) 的两个元素。存在唯一的元素 \(t\) 属于 \(\mathrm{D}_{\mathrm{ab}}^{*}\)，使得 \(\omega' = t\omega\)。

注. --- 假设体 D 是交换的。按定义，B(V) 是 \(V^{n}\) 的一个子集。显然，非零交错 n 重线性形式 \(f: V^{n} \to D\) 在 B(V) 上的限制属于 \(\Omega(V)\)。此外，若 \((e_{1}, \ldots, e_{n})\) 是 V 的一个基且 t 是 D 的一个非零元素，则存在唯一的交错 n 重线性形式 f，使得 \(f(e_{1}, \ldots, e_{n}) = t\)（III, §7, No.~4, 第 511 页与 III, §7, No.~8, 第 518 页）。由推论 1，集合 \(\Omega(V)\) 由诸非零交错 n 重线性形式在 B(V) 上的限制组成。

\section{自同构的行列式}

我们首先假设体 D 是交换的。鉴于上面的注，V 的自同构 u 的行列式是唯一的元素 \(\det u\) 属于 \(D^{*}\)，使得我们有

\[
\omega (u (v _ {1}), \ldots , u (v _ {n})) = (\det u) \omega (v _ {1}, \ldots , v _ {n})\tag{14}
\]

对 V 的每个基 \((v_{1},\ldots ,v_{n})\) 与每个元素 \(\omega\)（属于 \(\Omega (\mathrm{V})\)）成立。

现在回到不再假设 D 交换的情形。

命题 2. --- a) 设 u 是 V 的一个自同构。存在 \(D_{ab}^{*}\) 的唯一元素，记作 det u 并称为 u 的行列式，使得我们有

\[
\omega (u (v _ {1}), \ldots , u (v _ {n})) = (\det u) \omega (v _ {1}, \ldots , v _ {n})\tag{15}
\]

对 V 的每个基 \((v_{1},\ldots ,v_{n})\) 与每个 \(\omega\)（在 \(\Omega (\mathrm{V})\) 中）成立

\begin{enumerate}
\def\labelenumi{\alph{enumi})}
\setcounter{enumi}{1}
\tightlist
\item
  映射 \(u \mapsto \det u\)（从 \(\mathbf{GL}(\mathrm{V})\) 到 \(\mathrm{D}_{\mathrm{ab}}^{*}\)）是一个群同态。
\end{enumerate}

设 \(\omega_0\) 是 \(\Omega (\mathrm{V})\) 的一个元素。映射 \((v_{1},\ldots ,v_{n})\mapsto\) \(\omega_0(u(v_1),\dots ,u(v_n))\)（从 B(V) 到 \(\mathrm{D}_{\mathrm{ab}}^{*}\)）属于 \(\Omega (\mathrm{V})\)。由 VIII, 第 451 页的推论 2，存在唯一的元素 \(t\) 属于 \(\mathrm{D}_{\mathrm{ab}}^{*}\)，使得我们有

\[
\omega_ {0} (u (v _ {1}), \ldots , u (v _ {n})) = t \omega_ {0} (v _ {1}, \ldots , v _ {n})
\]

对 \((v_{1},\ldots ,v_{n})\)（在 \(\mathrm{B(V)}\) 中）成立。若 \(\omega\) 是 \(\Omega (\mathrm{V})\) 的另一个元素，则存在一个元素 \(s\) 属于 \(\mathrm{D}_{\mathrm{ab}}^{*}\)，使得

\[
\omega (v _ {1}, \ldots , v _ {n}) = s \omega_ {0} (v _ {1}, \ldots , v _ {n})
\]

对 V 的每个基 \((v_{1},\ldots ,v_{n})\) 成立（同前所引）。我们立即推出关系式

\[
\omega (u (v _ {1}), \dots , u (v _ {n})) = t \omega (v _ {1}, \dots , v _ {n}).
\]

这就证明了 a)；断言 b) 是 a) 的直接推论。

\section{方阵的行列式}

设 n 是一个正整数。我们把上述结果应用于体 D 上的右向量空间 \(D_{d}^{n}\)；\(D_{d}^{n}\) 的元素被视为 n 行一列的矩阵。设 \((\varepsilon_{1},\ldots,\varepsilon_{n})\) 是 \(D_{d}^{n}\) 的典范基。由 VIII, 第 451 页的推论 2，存在唯一的元素 \(\omega_{0}\) 属于 \(\Omega(\mathrm{D}_{d}^{n})\)，使得 \(\omega_{0}(\varepsilon_{1},\ldots,\varepsilon_{n})=1\)。若 A 是 \(\mathbf{GL}_{n}(\mathrm{D})\) 的一个元素，则它的列 \(a_{1},\ldots,a_{n}\) 构成 \(D_{d}^{n}\) 的一个基；元素 \(\omega_{0}(a_{1},\ldots,a_{n})\)（\(D_{ab}^{*}\) 中的）称为 A 的行列式，记作 \(\det(A)\)。由于我们有 \(a_{i}=A\varepsilon_{i}\)（对 \(1\leqslant i\leqslant n\)），A 的行列式就是自同构 \(x\mapsto Ax\)（\(D_{d}^{n}\) 的）的行列式。特别地，若体 D 交换，则 A 的行列式与 III, §8, No.~3, 第 524 页中定义的一致。

设 \(\mathrm{V}\) 是一个维数 \(n\) 有限的右向量空间（体 \(\mathrm{D}\) 上的），且 \((e_1, \ldots, e_n)\) 是 \(\mathrm{V}\) 的一个基。若 \(u\) 是 \(\mathrm{V}\) 的一个自同构，且 \(A\) 是 \(u\) 关于基 \((e_1, \ldots, e_n)\) 的矩阵，则我们有 \(\det(u) = \det(A)\)。

用 \((E_{ij})\) 表示 \(\mathbf{M}_{n}(\mathrm{D})\) 的矩阵单位族（II, §10, No.~3, 第 341 页）。对 D 的每个元素 \(\lambda\) 与每对不同的整数 \((i,j)\)（在 {[}1, n{]} 中），我们置（II, §10, No.~13, 第 361 页）

\[
B _ {i j} (\lambda) = I _ {n} + \lambda E _ {i j};
\]

它是 \(\mathbf{GL}_{n}(\mathrm{D})\) 的一个元素。若 \(\lambda_{1},\ldots,\lambda_{n}\) 是 \(D^{*}\) 的元素，则对角矩阵 \(\mathrm{diag}(\lambda_{1},\ldots,\lambda_{n})\) 也属于 \(\mathbf{GL}_{n}(\mathrm{D})\)。

命题 3. --- 映射 det 是从 \(\mathbf{GL}_{n}(\mathrm{D})\) 到 \(D_{ab}^{*}\) 的满足关系式

\[
\det (B _ {i j} (1)) = 1\tag{16}
\]

（对 \(i\neq j\)）与

\[
\det (\mathrm{diag} (\lambda_ {1}, \ldots , \lambda_ {n})) = \pi (\lambda_ {1} \dots \lambda_ {n}),\tag{17}
\]

（对 \(\lambda_1, \ldots, \lambda_n \in \mathrm{D}^*\)）的唯一同态。

设 \(A\) 是 \(\mathbf{GL}_n(\mathrm{D})\) 的一个元素，\(a_1, \ldots, a_n\) 是它的列。我们有

\[
\det (A) = \omega_ {0} (a _ {1}, \dots , a _ {n}).
\]

现在，矩阵 \(A \operatorname{diag}(\lambda_1, \ldots, \lambda_n)\) 的列是 \(a_1\lambda_1, \ldots, a_n\lambda_n\)，而矩阵 \(AB_{ij}(1)\) 的列因而是 \(a_1, \ldots, a_j + a_i, a_{j+1}, \ldots, a_n\)。由于 \(\omega_0\) 是 \(\Omega(\mathrm{D}_d^n)\) 中使 \(\omega_0(\varepsilon_1, \ldots, \varepsilon_n) = 1\) 的唯一元素，我们看出，行列式是满足下列关系式的唯一映射 \(\varphi: \mathbf{GL}_n(\mathrm{D}) \to \mathrm{D}_{\mathrm{ab}}^*\)

\begin{enumerate}
\def\labelenumi{(\arabic{enumi})}
\setcounter{enumi}{17}
\tightlist
\item
\end{enumerate}

\[
\varphi (A \operatorname{diag} (\lambda_ {1}, \ldots , \lambda_ {n})) = \varphi (A) \pi (\lambda_ {1} \dots \lambda_ {n}),\tag{19}
\]

\[
\varphi (A B _ {i j} (1)) = \varphi (A) \quad \mathrm{for} i \neq j,\tag{20}
\]

\[
\varphi (I _ {n}) = 1.
\]

这首先表明行列式映射满足关系式 (16) 与 (17)。我们已经知道这个映射是从 \(\mathbf{GL}_{n}(\mathrm{D})\) 到 \(D_{ab}^{*}\) 的一个同态。反之，若 \(\varphi\) 是从 \(\mathbf{GL}_{n}(\mathrm{D})\) 到 \(D_{ab}^{*}\) 的一个同态，使得 \(\varphi(B_{ij}(1)) = 1\)（对 \(i \neq j\)），且 \(\varphi(\text{diag}(\lambda_1, \ldots, \lambda_n)) = \pi(\lambda_1 \cdots \lambda_n)\)，则 \(\varphi\) 满足关系式 (18) 至 (20)，从而等于 det。

例. --- 1) 矩阵 \(B_{ij}(\lambda)\) 的列是 \(\varepsilon_1, \ldots, \varepsilon_j + \varepsilon_i \lambda, \varepsilon_{i+1}, \ldots, \varepsilon_n\)。由 VIII, 第 448 页的公式 (5)，我们有

\[
\det (B _ {i j} (\lambda)) = 1.\tag{21}
\]

\begin{enumerate}
\def\labelenumi{\arabic{enumi})}
\setcounter{enumi}{1}
\tightlist
\item
  设 \(\sigma\) 是区间 \([1,n]\)（在 \(\mathbf{N}\) 中）的一个置换，其符号差为 \(\varepsilon(\sigma)\)。设 \(M(\sigma)\) 是置换 \(\sigma\) 的矩阵（II, §10, No.~7, 第 351 页）。矩阵 \(M(\sigma)\) 的列是 \(\varepsilon_{\sigma(1)},\ldots,\varepsilon_{\sigma(n)}\)。于是应用 VIII, 第 448 页的公式 (4)，我们有
\end{enumerate}

\[
\det (M (\sigma)) = \pi (\varepsilon (\sigma)).\tag{22}
\]

\begin{enumerate}
\def\labelenumi{\arabic{enumi})}
\setcounter{enumi}{2}
\item
  设 \(n \geqslant 1\)。对每个形如 \(\Delta = \text{diag}(d_1, \ldots, d_n)\) 的可逆对角矩阵，我们有 \(\Delta B_{ij}(\lambda)\Delta^{-1} = B_{ij}(d_i\lambda d_j^{-1})\)。设 A 是 \(\mathbf{GL}_n(\mathrm{D})\) 的一个元素。由 II, §10, No.~13, 第 362 页的推论 1 以及前一个公式，存在矩阵 P 与 \(\Delta\)（在 \(\mathbf{GL}_n(\mathrm{D})\) 中），使得 \(A = P\Delta\)，P 是形如 \(B_{ij}(\lambda)\) 的矩阵的乘积，且 \(\Delta\) 是形如 \(\text{diag}(1, \ldots, 1, d)\) 的对角矩阵。由例 1 我们有 \(\det(P) = 1\)，从而由命题 3 有 \(\det(A) = \det(\Delta) = \pi(d)\)。
\item
  设 \(D'\) 是一个体，u 是从 D 到 \(D'\) 的一个同态。过渡到商群，u 定义一个群同态 \(u_{ab}\)（从 \(D_{ab}^{*}\) 到 \(D_{ab}^{\prime *}\)）。设 \(u_{n}\) 是从 \(\mathbf{GL}_{n}(\mathrm{D})\) 到 \(\mathbf{GL}_{n}(\mathrm{D}')\) 的把矩阵 \(A = (a_{ij})\) 变为矩阵 \((u(a_{ij}))\) 的同态。公式
\end{enumerate}

\[
\det (u _ {n} (A)) = u _ {\mathrm{ab}} (\det (A))\tag{23}
\]

对 \(A \in \mathbf{GL}_n(\mathrm{D})\) 立即由例 3 得出。

\begin{enumerate}
\def\labelenumi{\arabic{enumi})}
\setcounter{enumi}{4}
\item
  我们有 \(\mathbf{GL}_1(\mathrm{D}) = \mathrm{D}^*\)，且命题 3 表明我们有 \(\det(a) = \pi(a)\)（\(a\) 在 \(\mathbf{GL}_1(\mathrm{D})\) 中）。
\item
  设 \(n = 2\)。设 \(A = \begin{pmatrix} a & b \\ c & d \end{pmatrix}\) 是 \(\mathbf{GL}_2(\mathrm{D})\) 的一个元素。元素 \(a\) 与 \(c\)（\(\mathrm{D}\) 中的）不全为零。我们将明显地给出 \(A\) 的行列式。
\end{enumerate}

若 \(a\) 不为零，则我们有

\[
A = \left( \begin{array}{c c} 1 & 0 \\ c a ^ {- 1} & 1 \end{array} \right) \left( \begin{array}{c c} a & 0 \\ 0 & d - c a ^ {- 1} b \end{array} \right) \left( \begin{array}{c c} 1 & a ^ {- 1} b \\ 0 & 1 \end{array} \right).\tag{24}
\]

因此，\(ad - aca^{-1}b \neq 0\)，且

\[
\det (A) = \pi (a d - a c a ^ {- 1} b).\tag{25}
\]

\begin{enumerate}
\def\labelenumi{\alph{enumi})}
\setcounter{enumi}{1}
\tightlist
\item
  若 \(a\) 为零，则我们有 \(c \neq 0\)，且
\end{enumerate}

\[
A = \left( \begin{array}{c c} 0 & b \\ c & d \end{array} \right) = \left( \begin{array}{c c} 0 & 1 \\ 1 & 0 \end{array} \right) \left( \begin{array}{c c} c & d \\ 0 & b \end{array} \right).\tag{26}
\]

因此，由 a) 与例 2，我们有 \(cb \neq 0\)，且

\[
\det (A) = \pi (- c b).\tag{27}
\]

\begin{enumerate}
\def\labelenumi{\arabic{enumi})}
\setcounter{enumi}{6}
\tightlist
\item
  设 \(D^{o}\) 是 D 的反体。映射 \(A \mapsto {}^{t}A\)（从 \(\mathbf{M}_{n}(\mathrm{D})\) 到 \(\mathbf{M}_{n}(\mathrm{D}^{\circ})\)）满足 \({}^{t}(AB) = {}^{t}B^{t}A\)。因此，对 \(\mathbf{GL}_{n}(\mathrm{D})\) 中的每个矩阵 A，转置矩阵 \({}^{t}A\) 在 \(\mathbf{M}_{n}(\mathrm{D}^{\circ})\) 中可逆，但它在 \(\mathbf{M}_{n}(\mathrm{D})\) 中不一定可逆。由例 3 得出，若 A 属于 \(\mathbf{GL}_{n}(\mathrm{D})\)，则 \(\mathbf{GL}_{n}(\mathrm{D})\) 的元素 A 与 \({}^{t}A\)（\(\mathbf{GL}_{n}(\mathrm{D}^{\circ})\) 的）有相同的行列式。另一方面，即使矩阵 \({}^{t}A\) 属于 \(\mathbf{GL}_{n}(\mathrm{D})\)，它在 \(\mathrm{GL}_{n}(\mathrm{D})\) 中的行列式也不一定等于 A 的行列式（参见例 6）。
\end{enumerate}

注. --- 1) 设 V 是体 D 上的一个右向量空间，维数 n 有限。我们把对偶空间 \(V^{*} = \operatorname{Hom}_{\mathrm{D}}(\mathrm{V}, \mathrm{D})\) 视为 D 的反体 \(D^{o}\) 上的一个右向量空间。设 u 是 V 的一个自同构，\(^{t}u\) 是 \(V^{*}\) 的作为 u 的转置的自同构。若 u 由一个矩阵 A（在 \(\mathbf{M}_{n}(\mathrm{D})\) 中，关于 V 的一个基 \((e_{1}, \ldots, e_{n})\)）表示，则自同构 \(^{t}u\) 由矩阵 \(^{t}A\)（在 \(\mathbf{M}_{n}(\mathrm{D}^{o})\) 中）表示，所关于的基 \((e_{1}^{*}, \ldots, e_{n}^{*})\) 是 \(V^{*}\) 的与 \((e_{1}, \ldots, e_{n})\) 对偶的基（II, §10, No.~4, 第 344 页，命题 3）。由例 7，\(^{t}u\) 的行列式等于 u 的行列式。

\begin{enumerate}
\def\labelenumi{\arabic{enumi})}
\setcounter{enumi}{1}
\tightlist
\item
  第 1 至第 4 小节的结果可以推广到 D 是一个局部环的情形（VIII, 第 459 页，习题 2）。
\end{enumerate}

\section{幺模群}

设 n 是一个正整数。我们用 \(\mathbf{SL}_{n}(\mathrm{D})\) 表示群同态 \(\det: \mathbf{GL}_{n}(\mathrm{D}) \to \mathrm{D}_{\mathrm{ab}}^{*}\) 的核，并称它为幺模群或特殊线性群（当体 D 交换时，参见 III, §8, No.~9, 第 537 页）。

定理 1. --- 设 \(n \geqslant 2\)。

\begin{enumerate}
\def\labelenumi{\alph{enumi})}
\item
  子群 \(\mathbf{SL}_n(\mathrm{D})\)（\(\mathbf{GL}_n(\mathrm{D})\) 的）由诸矩阵 \(B_{ij}(\lambda)\) 生成，其中 \(i\) 与 \(j\) 是区间 \([1,n]\) 中两个不同的整数，而 \(\lambda\) 跑遍 D。
\item
  设 \(n \geqslant 3\) 或 \(\operatorname{Card}(\mathrm{D}) \geqslant 3\)。\(\mathbf{GL}_n(\mathrm{D})\) 的导群等于 \(\mathbf{SL}_n(\mathrm{D})\)。
\item
  设 \(n \geqslant 3\) 或 \(\operatorname{Card}(\mathrm{D}) \geqslant 4\)。\(\mathbf{SL}_n(\mathrm{D})\) 的导群等于 \(\mathbf{SL}_n(\mathrm{D})\)。
\end{enumerate}

\begin{enumerate}
\def\labelenumi{\Alph{enumi})}
\tightlist
\item
  用 T 表示 \(\mathbf{GL}_{n}(\mathrm{D})\) 的由诸矩阵 \(B_{ij}(\lambda)\) 生成的子群。由 VIII, 第 454 页的例 1，我们有 \(\det(B_{ij}(\lambda)) = 1\)，从而 \(\mathrm{T} \subset \mathbf{SL}_{n}(\mathrm{D})\)。为证明这两个群相等，于是由同前所引的例 3，只需证明每个形如 \(\operatorname{diag}(1, \ldots, 1, d)\) 且 \(\pi(d) = 1\) 的矩阵属于 T。矩阵 \(\operatorname{diag}(1, \ldots, 1, d)\) 属于同态 \(U \mapsto \begin{pmatrix} I_{n-2} & 0 \\ 0 & U \end{pmatrix}\)（从 \(\mathbf{GL}_{2}(\mathrm{D})\) 到 \(\mathbf{GL}_{n}(\mathrm{D})\)）的像；因此只需考虑 n = 2 的情形。由于 \(\pi\) 的核是 \(D^{*}\) 的导群，我们可以假设 \(d = uvu^{-1}v^{-1}\)，其中 u、v 属于 \(D^{*}\)。于是我们的断言由等式
\end{enumerate}

\[
\left( \begin{array}{c c} 1 & 0 \\ 0 & d \end{array} \right) = \left( \begin{array}{c c} u ^ {- 1} & 0 \\ 0 & u \end{array} \right) \left( \begin{array}{c c} v ^ {- 1} & 0 \\ 0 & v \end{array} \right) \left( \begin{array}{c c} v u & 0 \\ 0 & u ^ {- 1} v ^ {- 1} \end{array} \right)\tag{28}
\]

与

\[
\left( \begin{array}{c c} s & 0 \\ 0 & s ^ {- 1} \end{array} \right) = \left( \begin{array}{c c} 1 & s \\ 0 & 1 \end{array} \right) \left( \begin{array}{c c} 1 & 0 \\ 1 - s ^ {- 1} & 1 \end{array} \right) \left( \begin{array}{c c} 1 & - 1 \\ 0 & 1 \end{array} \right) \left( \begin{array}{c c} 1 & 0 \\ 1 - s & 1 \end{array} \right)\tag{29}
\]

（对 \(s \in D^{*}\)）得出。

\begin{enumerate}
\def\labelenumi{\Alph{enumi})}
\setcounter{enumi}{1}
\tightlist
\item
  按构造，\(\mathbf{SL}_{n}(\mathrm{D})\) 是一个从 \(\mathbf{GL}_{n}(\mathrm{D})\) 到交换群的同态的核，所以它包含导群 \((\mathbf{GL}_{n}(\mathrm{D}),\mathbf{GL}_{n}(\mathrm{D}))\)（\(\mathbf{GL}_{n}(\mathrm{D})\) 的）。由上述，我们有
\end{enumerate}

\[
\mathbf {S L} _ {n} (\mathrm{D}) \supset (\mathbf {G L} _ {n} (\mathrm{D}), \mathbf {G L} _ {n} (\mathrm{D})) \supset (\mathbf {S L} _ {n} (\mathrm{D}), \mathbf {S L} _ {n} (\mathrm{D})).
\]

为证明 c)，只需证明诸矩阵 \(B_{ij}(\lambda)\) 是 \(\mathbf{SL}_n(\mathrm{D})\) 中的换位子。

设 \(n \geqslant 3\)。若 \(i, j, k\) 是区间 \([1, n]\) 中三个不同的整数，且 \(\mu, \nu\) 是 \(D\) 的元素，则我们有

\[
B _ {i j} (\mu \nu) = B _ {i k} (\mu) ^ {- 1} B _ {k j} (\nu) ^ {- 1} B _ {i k} (\mu) B _ {k j} (\nu).\tag{30}
\]

取 \(\mu = 1\) 与 \(\nu = \lambda\)，我们看出矩阵 \(B_{ij}(\lambda)\) 是 \(\mathbf{SL}_n(\mathrm{D})\) 的元素的换位子。

现在设 \(n = 2\)。设 \(u\) 与 \(v\) 是 \(D\) 的元素且 \(u \neq 0\)。我们有关系式

\[
\left( \begin{array}{c c} 1 & v - u v u \\ 0 & 1 \end{array} \right) = \left( \begin{array}{c c} u & 0 \\ 0 & u ^ {- 1} \end{array} \right) \left( \begin{array}{c c} 1 & - v \\ 0 & 1 \end{array} \right) \left( \begin{array}{c c} u ^ {- 1} & 0 \\ 0 & u \end{array} \right) \left( \begin{array}{c c} 1 & v \\ 0 & 1 \end{array} \right)\tag{31}
\]

与

\[
\left( \begin{array}{c c} 1 & 0 \\ v - u v u & 1 \end{array} \right) = \left( \begin{array}{c c} u ^ {- 1} & 0 \\ 0 & u \end{array} \right) \left( \begin{array}{c c} 1 & 0 \\ - v & 1 \end{array} \right) \left( \begin{array}{c c} u & 0 \\ 0 & u ^ {- 1} \end{array} \right) \left( \begin{array}{c c} 1 & 0 \\ v & 1 \end{array} \right).\tag{32}
\]

我们有 \(\det\left(\begin{matrix}u & 0 \\ 0 & u^{-1}\end{matrix}\right)=1\)，所以矩阵 \(B_{12}(v-uvu)\) 与 \(B_{21}(v-uvu)\) 是 \(\mathbf{SL}_{n}(\mathrm{D})\) 的元素的换位子。

设体 D 至少有四个元素。设 \(\lambda\) 是 D 的一个元素。若 \(\lambda\) 等于 0、1 或 -1，则在 \(D-\{0,1,-1\}\) 中任意选取一个元素 u；否则，置 \(u=\lambda\)。在这两种情形下，u 都是 D 的非零元素，它与 \(\lambda\) 交换，且我们有 \(u^{2}\neq1\)。置 \(v=\lambda(1-u^{2})^{-1}\)。我们有 uv=vu，从而 \(v-uvu=v(1-u^{2})=\lambda\)。于是由关系式 (31) 与 (32) 得出，矩阵 \(B_{12}(\lambda)\) 与 \(B_{21}(\lambda)\) 是 \(\mathbf{SL}_{n}(\mathrm{D})\) 中的换位子。这就完成了 c) 的证明。

\begin{enumerate}
\def\labelenumi{\Alph{enumi})}
\setcounter{enumi}{2}
\tightlist
\item
  只剩下证明当 D 有三个元素时 \(\mathbf{SL}_{2}(\mathrm{D})\) 是 \(\mathbf{GL}_{2}(\mathrm{D})\) 的导群。由 A)，群 \(\mathbf{SL}_{2}(\mathrm{D})\) 由矩阵 \(B_{12}(1)=(\begin{matrix}{1}&{1}\\ {0}&{1}\end{matrix})\) 与 \(B_{21}(1)=(\begin{matrix}{1}&{0}\\ {1}&{1}\end{matrix})\) 生成，且这些矩阵是 \(\mathbf{GL}_{2}(\mathrm{D})\) 的元素的换位子，因为我们有 \(B_{21}(1)={}^{t}B_{12}(1)\)，且
\end{enumerate}

\[
\left( \begin{array}{c c} 1 & 1 \\ 0 & 1 \end{array} \right) = \left( \begin{array}{c c} - 1 & 0 \\ 0 & 1 \end{array} \right) \left( \begin{array}{c c} 1 & 1 \\ 0 & 1 \end{array} \right) \left( \begin{array}{c c} - 1 & 0 \\ 0 & 1 \end{array} \right) ^ {- 1} \left( \begin{array}{c c} 1 & 1 \\ 0 & 1 \end{array} \right) ^ {- 1}.\tag{33}
\]

注. --- 若 D 是一个只有两个元素的体，则 \(\mathbf{GL}_{2}(\mathrm{D})\) 等于 \(\mathbf{SL}_{2}(\mathrm{D})\)，它是一个 6 阶群，其导群的阶为 3。若 D 是一个只有三个元素的体，则群 \(\mathbf{SL}_{2}(\mathrm{D})\) 的阶为 24，其导群的阶为 8。设 \(n \geqslant 2\)；除前两种情形外，\(\mathbf{SL}_{n}(\mathrm{D})\) 的每个异于 \(\mathbf{SL}_{n}(\mathrm{D})\) 的正规子群都包含在中心 \(\mathrm{Z}(\mathbf{SL}_{n}(\mathrm{D}))\)（\(\mathbf{SL}_{n}(\mathrm{D})\) 的）中（II, §10, 第 421 页，习题 13 与 14）。这时群 \(\mathbf{SL}_{n}(\mathrm{D})/\mathrm{Z}(\mathbf{SL}_{n}(\mathrm{D}))\) 是单群。

设 V 是体 D 上的一个右向量空间，维数 n 有限。我们用 \(\mathrm{SL}(V)\) 表示同态 \(\det: \mathbf{GL}(V) \to D_{ab}^{*}\) 的核，并称之为 V 的幺模群。选取 V 的一个基就能把这个群与 \(\mathbf{SL}_{n}(D)\) 等同起来。

我们称 V 的一个自同构 u 为平延，如果存在 V 的一个向量 v 与 V 上的一个线性型 \(\varphi\)，使得 \(\varphi(v)=0\) 且对 V 中的每个 x 有 \(u(x)=x+v\varphi(x)\)。当 \(n\geqslant2\) 时，u 是平延当且仅当存在 V 的一个基，u 关于这个基的矩阵形如 \(B_{ij}(\lambda)\)。定理 1 蕴含下面的推论。

推论. --- 设 V 是体 D 上的一个右向量空间，维数 \(n \geqslant 2\)。

\begin{enumerate}
\def\labelenumi{\alph{enumi})}
\item
  子群 \(\mathbf{SL}(\mathrm{V})\)（\(\mathbf{GL}(\mathrm{V})\) 的）由诸平延生成。
\item
  子群 \(\mathbf{SL}(\mathrm{V})\) 是 \(\mathbf{GL}(\mathrm{V})\) 的导群，除非 \(n = 2\) 且 D 只有两个元素。
\item
  除非 n = 2 且 D 只有两个或三个元素，群 \(\mathbf{SL}(V)\) 等于它的导群。
\end{enumerate}

\BourbakiAppendixExercises{习题}

\begin{enumerate}
\def\labelenumi{\arabic{enumi})}
\item
  给出一个体 D 上的可逆矩阵 A 的例子，它的转置不可逆；再给出 D 上的一个可逆矩阵 B 的例子，它的转置可逆但满足 \(\det^{t}B \neq \det B\)（考虑四元数体上的 2 阶方阵）。
\item
  设 D 是一个局部环；我们用 \(D^{*}\) 表示 D 的可逆元素所成的群，用 \(D_{ab}^{*}\) 表示 \(D^{*}\) 对其导群的商，用 \(\pi : D^{*} \to D_{ab}^{*}\) 表示典范同态。设 V 是一个维数 n 有限的自由右 D-模。我们用 \(\mathrm{B(V)}\) 表示 V 的基所成的集合，用 \(\Omega(\mathrm{V})\) 表示满足下列两个条件的映射 \(\omega : \mathrm{B(V)} \to \mathrm{D}_{\mathrm{ab}}^{*}\) 所成的集合：
\end{enumerate}

\begin{enumerate}
\def\labelenumi{(\roman{enumi})}
\item
  对 D* 中的 \(\lambda_1, \ldots, \lambda_n\) 与 \((v_1, \ldots, v_n) \in \mathrm{B}(\mathrm{V})\)，我们有 \(\omega(v_1 \lambda_1, \ldots, v_n \lambda_n) = \omega(v_1, \ldots, v_n)\pi(\lambda_1 \ldots \lambda_n)\)。
\item
  对 \(i \neq j\) 与 \(\lambda \in D\)，我们有 \(\omega(v_{1},\ldots,v_{i}+v_{j}\lambda,\ldots,v_{n}) = \omega(v_{1},\ldots,v_{n})\)。a) 设 e 是 V 的一个元素，W 是 V 的一个与 eD 互补的子模。设 \(\varphi \in \Omega(\mathrm{W})\)；证明存在唯一的元素 \(\omega\) 属于 \(\Omega(\mathrm{V})\)，满足
\end{enumerate}

\[
\omega (w _ {1}, \dots , w _ {n - 1}, e) = \varphi (w _ {1}, \dots , w _ {n - 1})
\]

对 W 的每个基 \((w_{1},\ldots ,w_{n - 1})\) 成立（把命题 1 的证明加以变通）。

\begin{enumerate}
\def\labelenumi{\alph{enumi})}
\setcounter{enumi}{1}
\item
  由此推出：给定两个元素 \(\omega\) 与 \(\omega'\)（属于 \(\Omega(\mathrm{V})\)），存在唯一的元素 t 属于 \(D_{ab}^{*}\)，使得我们有 \(\omega' = t\omega\)。
\item
  设 u 是 V 的一个自同构。证明存在 \(D_{ab}^{*}\) 的唯一元素，记作 \(\det u\)，使得我们有 \(\omega(u(v_{1}),\ldots,u(v_{n}))=(\det u)\omega(v_{1},\ldots,v_{n})\)（对 V 的每个基 \((v_{1},\ldots,v_{n})\) 与每个元素 \(\omega\)，后者属于 \(\Omega(\mathrm{V})\)）。映射 \(u\mapsto\det u\)（从 \(\mathbf{GL}(\mathrm{V})\) 到 \(D_{ab}^{*}\)）是一个群同态。
\end{enumerate}

\(d\) ) 取 \(\mathrm{V} = \mathrm{D}_d^n\)，于是 \(\mathbf{GL}(\mathrm{V})\) 与群 \(\mathbf{GL}_n(\mathrm{D})\)（\(n\) 阶可逆方阵所成的群）等同。证明映射 \(\operatorname{det}:\mathbf{GL}_n(\mathrm{D})\to \mathrm{D}_{\mathrm{ab}}^*\) 是唯一的在所有矩阵 \(B_{ij}(\lambda)\)（\(i\neq j\)，\(\lambda \in \mathrm{D}\)）上取值 1 且满足 \(\operatorname{det}(\operatorname{diag}(\lambda_1,\dots ,\lambda_n)) = \pi (\lambda_1\dots \lambda_n)\)（对 \(\lambda_1,\ldots ,\lambda_n\)，在 \(\mathrm{D}^*\) 中）的群同态。

\begin{enumerate}
\def\labelenumi{\alph{enumi})}
\setcounter{enumi}{4}
\tightlist
\item
  把第 4 目的例 2 至例 6 变通到这个框架中。
\end{enumerate}

¶ 3) 设 \(\mathrm{A} = \mathrm{A}_0 \oplus \mathrm{A}_1\) 是一个 \(\mathbf{Z}/2\mathbf{Z}\) 型的分次环。假设 \(\mathrm{A}_0\) 包含在 \(\mathrm{A}\) 的中心中，且 \(\mathrm{A}_1\) 的每个元素的平方为零。

\begin{enumerate}
\def\labelenumi{\alph{enumi})}
\item
  设 \(A_{1}^{2}\) 是 \(A_{0}\) 的由 \(A_{1}\) 中两个元素的乘积生成的 Z-子模；证明 \(A_{1}^{2}\) 是 \(A_{0}\) 的一个诣零理想，且 \(A_{1}^{2} \oplus A_{1}\) 是 A 的一个双边诣零理想。
\item
  设 \(p, q\) 是 \(\geqslant 0\) 的整数。用 \(A^{p|q}\) 表示自由分次右 A-模 \(A_d^p \oplus A_d(1)^q\)（II, §11, No.~2, 第 365 页，例 3），用 \(\mathbf{GL}(p|q, \mathrm{A})\) 表示 \(A^{p|q}\) 的（分次）（0 次）自同构所成的群。这样一个自同构在 \(A^{p|q}\) 的典范基下由一个矩阵 \(X = \begin{pmatrix} A & B \\ C & D \end{pmatrix}\) 表示，其中 \(A \in \mathbf{M}_p(\mathrm{A}_0)\)、\(B \in \mathbf{M}_{p,q}(\mathrm{A}_1)\)、\(C \in \mathbf{M}_{q,p}(\mathrm{A}_1)\)、\(D \in \mathbf{M}_q(\mathrm{A}_0)\)。证明矩阵 \(A\) 与 \(D\) 可逆（利用 \(a\) 与 VIII, 第 168 页的习题 16）。
\item
  对上述的 X，置 \(\operatorname{sdet} X = \det(A - BD^{-1}C) \det D^{-1}\)。证明 \(\operatorname{sdet} X\) 是 \(A_{0}\) 的一个可逆元素。
\end{enumerate}

\(d\) ) 证明每个矩阵 \(X\in \mathbf{GL}(p|q,\mathrm{A})\) 可以写成一个乘积

\[
X = \left( \begin{array}{c c} I & R \\ 0 & I \end{array} \right) \left( \begin{array}{c c} P & 0 \\ 0 & Q \end{array} \right) \left( \begin{array}{c c} I & 0 \\ S & I \end{array} \right).
\]

（取 \(R = BD^{-1}\)、\(P = A - BD^{-1}C\)、\(Q = D\)、\(S = D^{-1}C\)。）

\begin{enumerate}
\def\labelenumi{\alph{enumi})}
\setcounter{enumi}{4}
\tightlist
\item
  证明 sdet 是从 \(\mathbf{GL}(p|q,\mathrm{A})\) 到 \(A_{0}\) 的可逆元素所成的群的一个同态。（利用 d) 归结为证明等式 \(\operatorname{sdet}(XY)=\operatorname{sdet}X\operatorname{sdet}Y\)，当 X 形如 \(\binom{I}{S}\operatorname{sdet}(I)\)（其中 S 是一个初等矩阵）时。注意，对每个矩阵 \(R\in\mathbf{M}_{p,q}(\mathrm{A}_{1})\)，矩阵 RS（相应地 SR）的任意两个元素的乘积为零，并由此推出等式 \(\det(I-RS)=(\det(I-SR))^{-1}\)。
\item
  设 \(X = \begin{pmatrix} A & B \\ C & D \end{pmatrix}\) 是 \(\mathbf{GL}(p|q,\mathrm{A})\) 的一个元素；置 \(^{st}X = \begin{pmatrix} t_{A} & t_{C} \\ -t_{B} & t_{D} \end{pmatrix}\) 与 \(^{\pi}X = \begin{pmatrix} D & C \\ B & A \end{pmatrix}\)。设 \(Y \in \mathbf{GL}(p|q, A)\)；我们有 \(^{st}(XY) = ^{st}Y^{st}X\) 与 \(^{\pi}(XY) = ^{\pi}X^{\pi}Y\)。证明等式 \(\operatorname{sdet}X = \operatorname{sdet}^{st}X = (\operatorname{sdet}^{\pi}X)^{-1}\)（利用 d) 计算 \(\operatorname{sdet}^{\pi}X\)）。
\end{enumerate}

\begin{enumerate}
\def\labelenumi{\arabic{enumi})}
\setcounter{enumi}{3}
\tightlist
\item
  保留习题 3 的假设。设 \(\mathrm{M} = \mathrm{M}_0 \oplus \mathrm{M}_1\) 是一个有限生成的自由分次右 A-模；存在唯一的整数 \(p, q\)，使得 \(\mathrm{M}\) 同构于 \(\mathrm{A}^{p|q}\)。我们称 \(\mathrm{M}\) 的一个基是适配的，如果它的前 \(p\) 个向量是 0 次齐次的，后 \(q\) 个是 1 次齐次的。
\end{enumerate}

\begin{enumerate}
\def\labelenumi{\alph{enumi})}
\item
  设 u 是分次模 M 的一个自同构，\(M_{\mathcal{B}}(u)\) 是它关于 B 的一个适配基的矩阵。证明元素 sdet \(M_{\mathcal{B}}(u)\) 不依赖于 B 的选取。我们把它记作 sdet u。
\item
  用 \(u(1)\) 表示把自同态 u 视为分次模 M(1) 的自同构。若 B 是 M 的一个适配基，则交换 B 的前 p 个向量与后 q 个向量所得到的基 \(^{\pi}B\) 是 M(1) 的一个适配基，且我们有 \(M_{\pi\mathcal{B}}(u(1)) = ^{\pi}M_{\mathcal{B}}(u)\)（习题 3, f)）。由此推出我们有 \(\operatorname{sdet}(u(1)) = (\operatorname{sdet} u)^{-1}\)。
\item
  在每个分次左 A-模 N 上，我们通过置 \(ax = (-1)^{ij}xa\)（\(a \in A_i\)，\(x \in N_j\)）定义一个典范的分次右 A-模结构。特别地，我们把 M 的分次对偶 \(M^{*gr}\) 视为一个分次右 A-模，它同样同构于 \(A^{p|q}\)。设 u 是 M 的一个自同构。我们用 \({}^{st}u\) 表示 \(M^{*gr}\) 的由公式 \(\langle{}^{st}u(x^{*}), x\rangle = (-1)^{rs}\langle x^{*}, u(x)\rangle\)（\(x^{*} \in (\mathrm{M}^{*\mathrm{gr}})_r\)，\(x \in M_s\)）刻画的自同构。若 B 是 M 的一个适配基，则对偶基 \(B^*\) 是 \(M^{*gr}\) 的一个适配基，且我们有 \(M_{\mathcal{B}^*}({}^{st}u) = {}^{st}M_{\mathcal{B}}(u)\)（习题 3, f）。我们有 \(\operatorname{sdet}({}^{st}u) = \operatorname{sdet}(u)\)。
\end{enumerate}

\chapter{希尔伯特零点定理}

定理 1. --- 设 A 是一个整环，K 是它的分式域，L 是一个交换 K-代数。假设 A-代数 L 由有限个元素生成，且 L 是一个体。

\begin{enumerate}
\def\labelenumi{\alph{enumi})}
\item
  L 在 K 上的次数有限。
\item
  存在一个非零元素 \(a\)（属于 \(\mathrm{A}\)），使得 \(\mathrm{K}\) 等于 \(\mathrm{A}[a^{-1}]\)。
\end{enumerate}

设 \(S\) 是 A-代数 \(L\) 的一个生成子集。我们对 \(S\) 的基数用归纳法。

首先，假设 S 的元素在 K 上都是代数的。这时 L 在 K 上的次数有限（V, §3, No.~2, 第 18 页，定理 2）。设 \((e_i)_{i\in I}\) 是 L 在 K 上的一个基。存在一个非零元素 \(a\)（属于 A），使得 S 的元素、元素 1 以及元素 \(e_i e_j\)（\(i\)、\(j\) 属于 I）关于这个基的坐标都属于 \(\mathrm{A}[a^{-1}]\)。于是诸线性组合 \(\sum_{i\in I}a_i e_i\)（其中 \(a_i\in \mathrm{A}[a^{-1}]\)，对每个 \(i\)）所成的集合是 L 的一个子环。按构造它包含 A 与 S，从而等于 L。特别地，它包含 \(\mathrm{Ke}_1\)，故 K 等于 \(\mathrm{A}[a^{-1}]\)。

现在设一个元素 \(s\)（属于 \(S\)）在 \(K\) 上是超越的。用 \(E\) 表示环 \(A[s]\) 的分式域。\(A[s]\)-代数 \(L\) 由 \(S - \{s\}\) 生成。由归纳假设，存在一个非零多项式 \(P \in A[X]\)，使得 \(E\) 等于 \(A[s][P(s)^{-1}]\)。设 \(\overline{K}\) 是 \(K\) 的一个代数闭包（V, §4, No.~3, 第 23 页，定理 2）。由于体 \(\overline{K}\) 是无限的（V, §4, No.~1, 第 20 页，命题 3），存在一个元素 \(x\)（属于 \(\overline{K}\)），使得 \(P(x) \neq 0\)。设 \(\varphi : E \to \overline{K}\) 是从 K-代数 \(E = A[s][P(s)^{-1}]\) 到 K-代数 \(\overline{K}\) 的把 \(s\) 映为 \(x\) 的唯一同态。这是荒谬的，因为 \(E\) 是 \(K\) 的一个超越扩张，而 \(\overline{K}\) 是 \(K\) 的一个代数扩张。这就完成了定理的证明。

推论 1. --- 设 K 是一个交换域，A 是一个由有限个元素生成的交换 K-代数，m 是 A 的一个极大理想。a) A/m 在 K 上的次数有限。

\begin{enumerate}
\def\labelenumi{\alph{enumi})}
\setcounter{enumi}{1}
\tightlist
\item
  设 \(\Omega\) 是 K 的一个代数闭扩张。存在一个从 A 到 \(\Omega\) 的以 \(\mathfrak{m}\) 为核的 K-代数同态。
\end{enumerate}

K-代数 A/m 由有限个元素生成，且 A/m 是一个体。由定理 1，A/m 的次数有限；由此得断言 a)。K 的每个次数有限的扩张都同构于 \(\Omega\) 的一个子扩张（V, §4, No.~1, 第 20 页，定理 1）；这就给出 b)。

推论 2. --- 设 K 是一个交换域，n 是一个自然数，\((\mathrm{P}_{i})_{i\in\mathrm{I}}\) 是 \(K[X_{1},\ldots,X_{n}]\) 的元素的一个族，\(\Omega\) 是 K 的一个代数闭扩张。下列条件等价：

\begin{enumerate}
\def\labelenumi{(\roman{enumi})}
\item
  多项式 \(\mathrm{P}_i\) 在 \(\Omega^n\) 中没有公共零点。
\item
  存在一个有限支撑族 \((\mathrm{Q}_i)_{i\in \mathrm{I}}\)（由 \(\mathrm{K}[\mathrm{X}_1,\dots ,\mathrm{X}_n]\) 的元素组成），使得 \(\sum_{i\in \mathrm{I}}\mathrm{P}_i\mathrm{Q}_i = 1\)。
\end{enumerate}

用 A 表示环 \(K[X_{1},\ldots,X_{n}]\)，用 a 表示由多项式 \(P_{i}\) 生成的理想。条件 (i) 表示不存在从 A/a 到 \(\Omega\) 的 K-代数同态。若它被满足，则由推论 1，环 A/a 没有极大理想，从而是零环，且 1 属于 a。这证明了 (i) 蕴含 (ii)。(ii) \(\Rightarrow\) (i) 是明显的。

\chapter{有限秩自同态的迹}

\section{有限秩线性映射}

设 A 是一个环，E 与 F 是 A-模。我们用 \(\operatorname{Hom}_{\mathrm{A}}^{\mathrm{f}}(\mathrm{E},\mathrm{F})\) 表示从 E 到 F 的像包含在 F 的一个有限生成子模中的线性映射所成的集合。它是 \(\operatorname{Hom}_{\mathrm{A}}(\mathrm{E},\mathrm{F})\) 的一个子群。当 A 是一个体时，它是从 E 到 F 的有限秩线性映射所成的集合（II, §7, No.~4, 第 298 页，定义 3）。置 \(\operatorname{End}_{\mathrm{A}}^{\mathrm{f}}(\mathrm{E}) = \operatorname{Hom}_{\mathrm{A}}^{\mathrm{f}}(\mathrm{E},\mathrm{E})\)。

设 E、F、G 是 A-模，\(u : E \to F\) 与 \(v : F \to G\) 是 A-线性映射。若 \(u \in \operatorname{Hom}_{\mathrm{A}}^{\mathrm{f}}(\mathrm{E}, \mathrm{F})\) 或 \(v \in \operatorname{Hom}_{\mathrm{A}}^{\mathrm{f}}(\mathrm{F}, \mathrm{G})\)，则 \(v \circ u\) 属于 \(\operatorname{Hom}_{\mathrm{A}}^{\mathrm{f}}(\mathrm{E}, \mathrm{G})\)。

用 \(\theta\) 表示从 \(\mathrm{E}^* \otimes_{\mathrm{A}} \mathrm{F}\) 到 \(\mathrm{Hom}_{\mathrm{A}}(\mathrm{E}, \mathrm{F})\) 的典范群同态（II, §4, No.~2, 第 271 页）；它把元素 \(x^* \otimes y\)（属于 \(\mathrm{E}^* \otimes_{\mathrm{A}} \mathrm{F}\)）映为映射 \(x \mapsto \langle x, x^* \rangle y\)（从 \(\mathrm{E}\) 到 \(\mathrm{F}\)）。

引理 1. --- 假设 A-模 F 是投射的。群同态 \(\theta\) 是单射，且它的像是 \(\mathrm{Hom}_{\mathrm{A}}^{\mathrm{f}}(\mathrm{E},\mathrm{F})\)。

由同前所引的推论，同态 \(\theta\) 是单射。它的像包含在 \(\operatorname{Hom}_{\mathrm{A}}^{\mathrm{f}}(\mathrm{E},\mathrm{F})\) 中。设 \(u \in \operatorname{Hom}_{\mathrm{A}}^{\mathrm{f}}(\mathrm{E},\mathrm{F})\)；让我们证明 u 属于 \(\theta\) 的像。

首先，假设 A-模 F 是自由的。设 \((f_{i})_{i\in\mathrm{I}}\) 是 F 的一个基，\((f_{i}^{*})_{i\in\mathrm{I}}\) 是 \(F^{*}\) 的与 \((f_{i})_{i\in\mathrm{I}}\) 对偶的基。设 J 是 I 的一个有限子集，使得 u 的像包含在由诸 \(f_{j}\)（\(j\in J\)）生成的 F 的子模中。我们有

\[
u (x) = \sum_ {j \in \mathrm{J}} \langle u (x), f _ {j} ^ {*} \rangle f _ {j} = \sum_ {j \in \mathrm{J}} \langle x, ^ {t} u (f _ {j} ^ {*}) \rangle f _ {j}\tag{1}
\]

对每个 \(x\in \mathrm{E}\) 成立，从而

\[
u = \theta \left(\sum_ {j \in \mathrm{J}} ^ {t} u \left(f _ {j} ^ {*}\right) \otimes f _ {j}\right).\tag{2}
\]

在一般情形，存在一个自由 A-模 L 与同态 \(i: F \to L\) 和 \(p: L \to F\)，使得 \(p \circ i = 1_{F}\)。同态 \(i \circ u\) 属于 \(\operatorname{Hom}_{\mathrm{A}}^{\mathrm{f}}(\mathrm{E}, \mathrm{L})\)。由上述，存在一个有限集 J，且对每个 \(j \in J\)，存在元素 \(x_{j}^{*}\)（属于 \(E^{*}\)）与 \(y_{j}\)（属于 L），使得我们有

\[
i (u (x)) = \sum_ {j \in \mathrm{J}} \langle x, x _ {j} ^ {*} \rangle y _ {j}
\]

对每个 \(x \in \mathrm{E}\) 成立。于是我们有

\[
u (x) = p \left(i (u (x))\right) = \sum_ {j \in J} \langle x, x _ {j} ^ {*} \rangle p \left(y _ {j}\right)
\]

对每个 \(x \in \mathrm{E}\) 成立，从而 \(u = \theta\big(\sum_{j \in \mathrm{J}} x_j^* \otimes p(y_j)\big)\)。

\section{有限秩自同态的迹}

在本小节中，A 表示一个交换环。设 E 是一个投射 A-模。这时 \(\operatorname{End}_{\mathrm{A}}^{\mathrm{f}}(\mathrm{E})\) 是 \(\operatorname{End}_{\mathrm{A}}(\mathrm{E})\) 的一个 A-子模，且典范映射 \(E^{*}\otimes_{A}E\to\operatorname{End}_{\mathrm{A}}(\mathrm{E})\) 定义 A-模的一个同构 \(\theta_{E}\)（从 \(E^{*}\otimes_{A}E\) 到 \(\operatorname{End}_{\mathrm{A}}^{\mathrm{f}}(\mathrm{E})\)）（引理 1）。考虑典范线性型 \(\tau:E^{*}\otimes_{A}E\to A\)（II, §4, No.~3, 第 273 页），它由公式 \(\tau(x^{*}\otimes x)=\langle x,x^{*}\rangle\) 刻画。把它与同构 \(\theta_{E}^{-1}\) 复合，我们得到一个线性型 \(\operatorname{Tr}: \operatorname{End}_{\mathrm{A}}^{\mathrm{f}}(\mathrm{E})\to\mathrm{A}\)，称为迹型。当 A-模 E 有限生成时，我们重新得到 II, §4, No.~3, 第 273 页中的定义。

命题 1. --- 设 E 是一个自由 A-模。设 \((e_{i})_{i\in\mathrm{I}}\) 是 E 的一个基，\((e_{i}^{*})_{i\in\mathrm{I}}\) 是它的对偶基。设 \(u\in\operatorname{End}_{\mathrm{A}}^{\mathrm{f}}(\mathrm{E})\)。族 \((\langle u(e_{i}),e_{i}^{*}\rangle)_{i\in\mathrm{I}}\) 有有限支撑，且它的和等于 \(\operatorname{Tr}(u)\)。

只需处理 \(u\) 形如 \(\theta_{\mathrm{E}}(x^{*} \otimes x)\)（其中 \(x \in \mathrm{E}\) 且 \(x^{*} \in \mathrm{E}^{*}\)）的情形。这时族 \((\langle x, e_{i}^{*} \rangle)_{i \in \mathrm{I}}\) 有有限支撑，且我们有 \(x = \sum_{i \in \mathrm{I}} \langle x, e_{i}^{*} \rangle e_{i}\)。于是族 \((\langle x, e_{i}^{*} \rangle \langle e_{i}, x^{*} \rangle)_{i \in \mathrm{I}}\) 也有有限支撑，且我们有 \(\langle x, x^{*} \rangle = \sum_{i \in \mathrm{I}} \langle x, e_{i}^{*} \rangle \langle e_{i}, x^{*} \rangle\)。现在，我们有 \(\langle u(e_{i}), e_{i}^{*} \rangle = \langle x, e_{i}^{*} \rangle \langle e_{i}, x^{*} \rangle\)（对每个 \(i \in \mathrm{I}\)）。这就证明了该命题。

命题 2. --- 设 E、F 是投射 A-模。设 \(u \in \operatorname{Hom}_{\mathrm{A}}^{\mathrm{f}}(\mathrm{E}, \mathrm{F})\)，\(v \in \operatorname{Hom}_{\mathrm{A}}(\mathrm{F}, \mathrm{E})\)。我们有关系式

\[
\operatorname{Tr} (v \circ u) = \operatorname{Tr} (u \circ v).\tag{3}
\]

只需在 \(u\) 形如 \(\theta(x^{*} \otimes y)\)（其中 \(x^{*} \in \mathrm{E}^{*}\) 且 \(y \in \mathrm{F}\)）时证明该命题。在这种情形，我们有

\[
v \circ u = \theta_ {\mathrm{E}} (x ^ {*} \otimes v (y)) \quad \text { and } \quad u \circ v = \theta_ {\mathrm{F}} (^ {t} v (x ^ {*}) \otimes y)  ,
\]

从而

\[
\operatorname{Tr} (v \circ u) = \langle v (y), x ^ {*} \rangle = \langle y, ^ {t} v (x ^ {*}) \rangle = \operatorname{Tr} (u \circ v).
\]

推论. --- 设 E 是一个投射 A-模，u 是 \(\operatorname{End}_{\mathrm{A}}^{f}(\mathrm{E})\) 的一个元素，F 是 E 的一个包含 Im u 的投射 A-子模。用 \(u_{F}\) 表示由 u 诱导的 F 的自同态。我们有

\[
\operatorname{Tr} (u) = \operatorname{Tr} (u _ {\mathrm{F}}).\tag{4}
\]

用 i 表示 F 到 E 中的典范单射，用 \(v: E \rightarrow F\) 表示由 u 得到的同态。我们有 \(u_{F} = v \circ i\) 与 \(u = i \circ v\)。由此即得该推论。

设 E 是一个投射 A-模，\(u \in \operatorname{End}_{\mathrm{A}}^{\mathrm{f}}(\mathrm{E})\)。对每个自然数 p，A-模 \(\wedge^{p}E\) 是投射的（III, §7, No.~8, 第 519 页，推论 2），且自同态 \(\wedge^{p}u\) 属于 \(\operatorname{End}_{\mathrm{A}}^{\mathrm{f}}(\wedge^{p}\mathrm{E})\)（III, §7, No.~3, 第 511 页，命题 6），并对充分大的 p 为零。集合 \(1_{\mathrm{E}} + \operatorname{End}_{\mathrm{A}}^{\mathrm{f}}(\mathrm{E})\) 对复合是稳定的。我们定义一个从 \(1_{\mathrm{E}} + \operatorname{End}_{\mathrm{A}}^{\mathrm{f}}(\mathrm{E})\) 到 A 的映射 det，置

\[
\det (1 _ {\mathrm{E}} + u) = \sum_ {p \geqslant 0} \operatorname{Tr} \wedge^ {p} u
\]

（对 \(u\in \mathrm{End}_{\mathrm{A}}^{\mathrm{f}}(\mathrm{E})\)）。

若 E 是自由的且有限维，则由 III, §8, No.~5, 第 530 页的推论，这个定义与 III, §8, No.~1, 第 522 页中的定义一致。

命题 3. --- 设 E 是一个投射 A-模。

\begin{enumerate}
\def\labelenumi{\alph{enumi})}
\tightlist
\item
  设 \(u \in \mathrm{End}_{\mathrm{A}}^{\mathrm{f}}(\mathrm{E})\)。设 \(\mathrm{F}\) 是 \(\mathrm{E}\) 的一个包含 \(\operatorname{Im} u\) 的投射 A-子模，\(u_{\mathrm{F}}\) 是 \(\mathrm{F}\) 的由 \(u\) 诱导的自同态。我们有
\end{enumerate}

\[
\det (1 _ {\mathrm{E}} + u) = \det (1 _ {\mathrm{F}} + u _ {\mathrm{F}}).\tag{5}
\]

\begin{enumerate}
\def\labelenumi{\alph{enumi})}
\setcounter{enumi}{1}
\tightlist
\item
  设 \(u\) 与 \(v\) 是 \(\mathrm{End}_{\mathrm{A}}^{\mathrm{f}}(\mathrm{E})\) 的两个元素。我们有
\end{enumerate}

\[
\det \bigl ((1 _ {\mathrm{E}} + u) \circ (1 _ {\mathrm{E}} + v) \bigr) = \det (1 _ {\mathrm{E}} + u) \det (1 _ {\mathrm{E}} + v).\tag{6}
\]

让我们证明 a)。对每个整数 \(p \geqslant 0\)，投射 A-模 \(\wedge^p\mathrm{F}\) 可以与 \(\wedge^p\mathrm{E}\) 的一个子模等同（III, §7, No.~9, 第 520 页，推论）。\(\wedge^p u\) 的像包含在 \(\wedge^p\mathrm{F}\) 中，且 \(\wedge^p\mathrm{F}\) 的由 \(\wedge^{p}u\) 诱导的自同态等于 \(\wedge^{p}u_{F}\)。于是由命题 2 的推论，我们有 \(\operatorname{Tr}(\wedge^{p}u)=\operatorname{Tr}(\wedge^{p}u_{F})\)，从而得 a)。

让我们证明 b)。设 G 是一个使 A-模 \(L = E \oplus G\) 为自由的 A-模。用 \(u'\) 与 \(v'\) 表示 L 的自同态 \(u \oplus 0_{G}\) 与 \(v \oplus 0_{G}\)。由 a)，我们有关系式 \(\det(1_{L} + u') = \det(1_{E} + u)\)、\(\det(1_{L} + v') = \det(1_{E} + v)\)，以及

\[
\det (1 _ {\mathrm{L}} + u ^ {\prime} + v ^ {\prime} + u ^ {\prime} \circ v ^ {\prime}) = \det (1 _ {\mathrm{E}} + u + v + u \circ v).
\]

因此只需在 A-模 E 为自由时证明断言 b)。这时存在 E 的一个有限生成自由子模 F，它包含 u 的像与 v 的像。置 \(w = u + v + u \circ v\)。w 的像包含在 F 中，且我们有 \(w_{F} = u_{F} + v_{F} + u_{F} \circ v_{F}\)。于是由 (5)，我们有 \(\det(1_{E} + u) = \det(1_{F} + u_{F})\)、\(\det(1_{E} + v) = \det(1_{F} + v_{F})\)，以及

\[
\begin{array}{c} \det \bigl ((1 _ {\mathrm{E}} + u) \circ (1 _ {\mathrm{E}} + v) \bigr) \\ = \det (1 _ {\mathrm{E}} + w) = \det (1 _ {\mathrm{F}} + w _ {\mathrm{F}}) = \det \bigl ((1 _ {\mathrm{F}} + u _ {\mathrm{F}}) \circ (1 _ {\mathrm{F}} + v _ {\mathrm{F}}) \bigr). \end{array}
\]

由于 F 是一个有限生成的自由 A-模，我们有

\[
\det \left(\left(1 _ {\mathrm{F}} + u _ {\mathrm{F}}\right) \circ \left(1 _ {\mathrm{F}} + v _ {\mathrm{F}}\right)\right) = \det \left(1 _ {\mathrm{F}} + u _ {\mathrm{F}}\right) \det \left(1 _ {\mathrm{F}} + v _ {\mathrm{F}}\right) = \det \left(1 _ {\mathrm{E}} + u\right) \det \left(1 _ {\mathrm{E}} + v\right).
\]

\BourbakiAppendixExercises{习题}

\begin{enumerate}
\def\labelenumi{\arabic{enumi})}
\tightlist
\item
  设 K 是一个交换域，E 是 K 上的一个向量空间。我们用 \(\operatorname{End}_{\mathrm{K}}^{\mathrm{pf}}(\mathrm{E})\) 表示 \(\operatorname{End}_{\mathrm{K}}(\mathrm{E})\) 中由某个幂属于 \(\operatorname{End}_{\mathrm{K}}^{\mathrm{f}}(\mathrm{E})\) 的自同态组成的子集。若 \(u \in \operatorname{End}_{\mathrm{K}}^{\mathrm{pf}}(\mathrm{E})\)，则置 \(I_{u} = \bigcap_{n \geqslant 0} Im u^{n}\) 与 \(\operatorname{Tr} u = \operatorname{Tr}(u | I_{u})\)；对有限秩自同态，这个概念与第 2 目中定义的一致。a) 设 \(u \in \operatorname{End}_{\mathrm{K}}^{\mathrm{pf}}(\mathrm{E})\)；设 V 是 E 的一个在 u 下稳定的子空间，\(u_{V}\) 与 \(u_{E/V}\) 是由 u 诱导的 V 与 E/V 的自同态。证明等式 \(\operatorname{Tr} u = \operatorname{Tr} u_{V} + \operatorname{Tr} u_{E/V}\)。
\end{enumerate}

\begin{enumerate}
\def\labelenumi{\alph{enumi})}
\setcounter{enumi}{1}
\item
  设 U 是 \(\operatorname{End}_{\mathbf{K}}(\operatorname{E})\) 的一个子空间。假设存在一个整数 n，使得 U 中 n 个元素的每个乘积都有有限秩。证明映射 \(Tr: U \to K\) 是线性的。
\item
  设 F 是 K 上的一个向量空间，\(u: E \to F\) 与 \(v: F \to E\) 是使得 \(u \circ v \in \operatorname{End}_{\mathrm{K}}^{\mathrm{pf}}(\mathrm{F})\) 的同态。我们有 \(v \circ u \in \operatorname{End}_{\mathrm{K}}^{\mathrm{pf}}(\mathrm{E})\) 与 \(\operatorname{Tr} v \circ u = \operatorname{Tr} u \circ v\)。
\end{enumerate}

¶ 2) 保留上一个习题的记号。若 V 与 W 是 E 的子空间，则我们用 V \(\prec\) W 表示关系“W 在 V + W 中具有有限余维数”。

\begin{enumerate}
\def\labelenumi{\alph{enumi})}
\item
  设 V 是 E 的一个子空间。我们用 \(\mathfrak{A}\)（相应地 \(\mathfrak{a},\mathfrak{b}\)）表示 V 的自同态 \(u\) 中满足 \(u(\mathrm{V})\prec \mathrm{V}\)（相应地 \(u(\mathrm{E})\prec \mathrm{V}\)、\(u(\mathrm{V})\prec 0\)）者所成的集合。证明 \(\mathfrak{A}\) 是 \(\operatorname{End}_{\mathrm{K}}(\mathrm{E})\) 的一个 K-子代数，\(\mathfrak{a}\) 与 \(\mathfrak{b}\) 是 \(\mathfrak{A}\) 的两个双边理想，其和为 \(\mathfrak{A}\)，且我们有 \(\mathfrak{a} \cap \mathfrak{b} \subset \operatorname{End}_{\mathrm{K}}^{\mathrm{pf}}(\mathrm{E})\)。
\item
  设 \(u, u'\) 是 \(\mathfrak{A}\) 中可交换的元素；我们可以写 \(u = a + b\) 与 \(u' = a' + b'\)，其中 \(a, a'\) 属于 \(\mathfrak{a}\)，\(b, b'\) 属于 \(\mathfrak{b}\)。证明元素 \([a, a'] = aa' - a'a\)（\(\mathfrak{A}\) 中的）属于 \(\mathfrak{a} \cap \mathfrak{b}\)，且它的迹不依赖于 \(u\) 与 \(u'\) 的分解的选取；我们把它记作 \(\langle u, u' \rangle\)。
\item
  设 A 是 \(\mathfrak{A}\) 的一个交换子代数。证明存在唯一的 K-线性映射 \(\operatorname{Res}:\Omega_{\mathrm{K}}(\mathrm{A})\to \mathrm{K}\)，使得 \(\operatorname{Res}(udv) = \langle u,v\rangle\)（对 A 中所有的 \(u,v\)）（利用恒等式 \([u,vw] + [v,wu] + [w,uv] = 0\)）。
\item
  设 \(u, v \in A\)；假设 \(v(V) \subset V\)。设 \(\pi\) 是 E 的以 V 为像的投影。证明等式 \(\operatorname{Res}(u dv) = \operatorname{Tr}[\pi u, v]\)。由此推出，若 \(u(V) \subset V\)，则我们有 \(\operatorname{Res} u dv = 0\)（注意自同态 \([\pi u, v]\) 的平方为零）。
\item
  设 \(u \in \mathrm{A}\)，\(n \in \mathbf{N}\)。证明我们有 \(\operatorname{Res} u^n du = 0\)。若 \(u\) 在 \(\mathrm{A}\) 中可逆，则我们有 \(\operatorname{Res} u^{-n} du = 0\)（对 \(n \geqslant 2\)）；若此外 \(u(\mathrm{V}) \subset \mathrm{V}\)，则我们有 \(\operatorname{Res} u^{-1} du = \dim_{\mathrm{K}} \mathrm{V} / u(\mathrm{V})\)。
\end{enumerate}

\(f)\) 取 \(\mathrm{E} = \mathrm{K}((\mathrm{X}))\)（IV, §4, No.~9, 第 38 页），\(\mathrm{V} = \mathrm{K}[[\mathrm{X}]]\)，并让 \(\mathrm{A} = \mathrm{E}\) 通过位似作用在自身上。证明 \(c)\) 的假设被满足。设 \(f(\mathrm{X}) = \sum_{n} a_{n} \mathrm{X}^{n} d\mathrm{X}\) 是 \(\mathrm{K}((\mathrm{X}))\) 的一个元素。证明 \(\operatorname{Res} f(\mathrm{X}) d\mathrm{X}\) 等于 \(a_{-1}\)。由此推出：若 \(g\) 是 \(\mathrm{K}[[\mathrm{X}]]\) 的一个元素，满足 \(g(0) = 0\) 且 \(g'(0) \neq 0\)，则 \(\mathrm{X}^{-1}\) 在 \(f(\mathrm{X})\) 与 \(f(g(\mathrm{X})) g'( \mathrm{X})\) 的形式级数展开中的系数相同。

\backmatter
\BourbakiBackMatter{历史注记}

我们已经看到（第一章与第二、三章的历史注记），最初的非交换代数出现于 1843--44 年，见于 Hamilton {[}25{]} 与 Grassmann {[}23{]}、{[}24{]} 的著作中。Hamilton 在引入四元数时，就已经对实数域上的任意有限秩代数有了清晰的想法（{[}25{]}，序言，第 26--31 页）\(^{(1)}\)。在发展他的理论的过程中，他后来产生了考虑他所谓的“双四元数”的想法，即复数域上与四元数代数有相同乘法表的代数；在这个场合，他注意到这种扩张导致零因子的出现（{[}25{]}，第 650 页）。Grassmann 的观点略有不同，在很长一段时间里，他的“外代数”一直与一般代数理论相互分离；\(^{(2)}\) 但在他那尚欠精确的语言中，人们不能不辨认出由生成元组与关系定义的（实数域上的、有限维或非有限维的）代数的最初思想 {[}24{]}。

新的代数例子在 1850 年至 1860 年间被或多或少明确地引入。Cayley 在发展矩阵理论时 {[}8{]}，尚未把方阵视为构成一个代数（这一观点最先由 Peirce 父子于 1870 年前后清楚地表达 {[}40{]}）。不过，他已在这个场合注意到存在一个满足四元数乘法表的 2 阶矩阵组；这一观察可以看作代数的线性表示的第一个例子。\(^{(3)}\) 此外，在他定义有限群的抽象概念的那篇论文中，他还顺便给出了这样一个群的代数的定义，但没有进一步应用这个定义（{[}7{]}，第 129 页）。

1870 年之前没有别的进展，但正是在这个时候，对（实数域或复数域上的）有限维代数一般结构的研究开始了。B. Peirce 在这个方向上迈出了最初的步伐；他引入了幂零元素与幂等元素的概念，证明了一个（有或没有单位元的）代数只要含有一个非幂零元素就有一个非零幂等元，写下了著名的分解

\[
x = e x e + (x e - e x e) + (e x - e x e) + (x - x e - e x + e x e)
\]

（e 为一个幂等元，x 为任意元素），并且有了把一个幂等元分解为两两正交的“本原”幂等元之和的（仍略欠精确的）想法 {[}40{]}。此外，据 Clifford 所说（{[}11{]}，第 274 页）\(^{(4)}\)，两个代数的张量积这一概念应归功于 B. Peirce；Clifford 本人在若干年后把它隐式地应用于 Hamilton“双四元数”的一个推广 {[}9{]}，并显式地应用于以他的名字命名的代数的研究（{[}10{]} 与 {[}11{]}）。这些新概念被 B. Peirce 用来对（复数域上的）低维代数进行分类。这个问题在大约 1880 年也被英美学派的其他数学家处理过，特别是 Cayley 与 Sylvester。可能的结构五花八门这一事实很快显露出来，无疑正是这一事实使随后时期的研究转向获取具有更特殊性质的代数类。

在欧洲大陆，思想演进的方式颇为不同，这类研究在 1880 年之前即已出现。1878 年，Frobenius 证明四元数构成实数域上唯一的非交换（有限维）体的例子（{[}18{]}，第 59--63 页），这一结果两年后由 C. S. Peirce 独立发表 {[}39{]}。早在 1861 年，Weierstrass 在澄清 Gauss 的一个评注时，就在他的讲课中把 R 或 C 上不含幂零元素的交换代数 \(^{(5)}\) 刻画为同构于 R 或 C 的体的直和。Dedekind 在他的研究中，于 1870 年前后结合其交换域理论的“超复”观念得到了同样的结论。他们的证明发表于 1884--85 年（{[}54{]} 与 {[}12{]}）。也是在 1884 年，H. Poincaré 在一篇十分隐晦的注记 {[}41{]} 中，提请注意这样一种可能性：把方程组 \(z_{i} = \varphi_{i}(x_{1}, \ldots, x_{n}, y_{1}, \ldots, y_{n})\)（它表达一个代数中的乘法法则 \((\sum_{i} x_{i} e_{i})(\sum_{i} y_{i} e_{i}) = \sum_{i} z_{i} e_{i}\)）视为（当然是在局部）定义一个李群。这个评注似乎给 Lie 及其追随者们（E. Study、G. Scheffers、F. Schur，稍后还有 T. Molien 与 É. Cartan）留下了深刻印象，他们当时正在发展“连续”群理论，尤其是分类问题（特别参见 {[}43{]}，第 387 页）。在 1885--1905 年间，这促使该学派的数学家把与研究李群和李代数时所用的类似方法应用于代数结构的研究。这些方法首先基于考虑代数的元素关于其正则表示的特征多项式（这个多项式已见于上面提到的 Weierstrass 与 Dedekind 的著作中），以及把这个多项式分解为不可约因子。正如 Frobenius 稍后发现的那样，这个分解反映在正则表示分解为不可约分量之中。

在 Lie 学派关于代数的研究过程中，这门理论的“内在”概念逐渐显现。根的概念于 1891 年在 G. Scheffers 的著作 {[}43{]} 中以一种特殊情形（即对根的商是体的直积的情形）出现，并在研究一般情形的 T. Molien {[}31{]} 与 É. Cartan {[}4{]} 的著作中更为清晰（“根”一词来自 Frobenius {[}22{]}）。Study 与 Scheffers {[}43{]} 突出了一个代数是若干个其他代数的直积的概念（Peirce 已有萌芽（{[}38{]}，第 221 页））。最后，Molien {[}31{]} 引入了一个代数的商代数，这一概念本质上等价于双边理想（最先由 É. Cartan 定义 {[}4{]}）或同态（这个名称也来自 Frobenius）。与群论的类比是显然的，后来在 1904 年，Epsteen 与 Wedderburn 考虑了双边理想的合成列并把若尔当--赫尔德定理推广到它们上。这个时期最重要的结果是 T. Molien 的 {[}31{]}：在单群概念的指引下，他定义了（C 上的）单代数并证明这些代数是矩阵代数。接着他证明，C 上任意有限秩代数的结构本质上归结为由 Scheffers 研究过的对根的商是体的直和的情形。此后不久，这些结果被 É. Cartan {[}4{]} 重新发现，并建立得更严格、更清晰。在这个场合，他引入了半单代数的概念，并揭示了与 C 上任意代数相伴的数值不变量，即“É. Cartan 整数”（VIII, 第 180 页，习题 8）。他就这样把这门代数的理论带到了一个此后鲜有进展的地步。\(^{(6)}\) 最后，他把 Molien 的以及他自己关于 R 上代数的结果加以推广。

大约在 1900 年，思想开始朝这样一个方向发展：在与线性代数有关的一切事情中放弃对标量体的任何限制。特别地，我们必须指出以 E. H. Moore 与 L. E. Dickson 为中心的美国学派对有限域研究所给予的强劲推动。这项研究的成果是 Wedderburn 定理 {[}50{]}，它证明每个有限域都是交换的。1907 年，Wedderburn 接过 Cartan 的结果并把它们推广到任意基域上 {[}51{]}。在这样做时，他完全放弃了其前辈们的方法（当体不再是代数闭的或极大序的时候，那些方法不再适用）。他转而对 B. Peirce 的幂等元技巧加以提炼后重新使用，这使他得以把半单代数的结构定理化为确定的形式，把半单代数的研究归结为非交换体的研究。此外，标量扩张的问题在他的视角下自然出现，Wedderburn 证明每个半单代数在基域的可分扩张后仍是半单的 \(^{(7)}\)，且当这个扩张足够大时成为中心单矩阵代数的直积（{[}51{]}，第 102 页）。\(^{(8)}\) 稍后，Dickson 对 n = 3 {[}17{]}，以及 Wedderburn 本人对任意 n {[}52{]}，给出了在其中心上秩为 \(n^{2}\) 的非交换体的最初例子，\(^{(9)}\) 从而在特殊情形引入了后来由 R. Brauer {[}2{]} 与 E. Noether（{[}33{]}、{[}34{]} 与 {[}35{]}）发展的“交叉乘积”与“因子系”理论。最后，在 1921 年，Wedderburn 证明了交换子定理的一个特殊情形 {[}53{]}。

与此同时，从 1896 年到 1910 年，一门与代数理论相近的理论，即群的线性表示理论（起初仅限于有限群的表示），由 Frobenius、Burnside 与 I. Schur 发展起来。它源于 Dedekind 所作的一些评注。在研究伽罗瓦扩张的正规基的过程中，Dedekind（甚至还在他关于代数的著作发表之前）大约在 1880 年遇到了“群行列式” \(\det(x_{st^{-1}})\)，其中 \((x_{s})_{s\in G}\) 是以有限群 G 为指标集的一个变量序列（换句话说，就是群 G 的代数的泛元素关于其正则表示的范数）。他观察到，当 G 交换时，这个多项式分解为线性因子（这推广了人们在很久以前对相应于循环群 G 的“循环”矩阵的行列式所证明的一个恒等式）。在他与 Frobenius 的引人入胜的通信 {[}13{]} 中，Dedekind 于 1896 年使后者注意到这一性质、它与交换群特征标理论的联系，以及他于 1886 年得到的关于某些特定非交换群的类似结果。几个月后，Frobenius 借助他对特征标概念的出色推广 {[}20{]}（我们不在此讨论它）完全解决了“群行列式”分解为不可约因子的问题 {[}19{]}。然而应当指出，在这一理论后来的发展中，\(^{(10)}\) Frobenius 始终意识到它与代数理论的联系（Dedekind 在他的信中从未停止强调这一点）。在引入了群的不可约表示与完全可约表示的概念 {[}21{]} 并证明了正则表示包含所有的不可约表示之后，正是通过类似的方法，他于 1903 年着手重新研究 Molien--Cartan 的理论 {[}22{]}。在 Burnside {[}3{]} 与 I. Schur {[}45{]} 的工作中，这一理论的“超复”方面并未明显地介入；但正是在他们的工作中，出现了不可约表示的基本性质：舒尔引理与伯恩赛德定理。最后还应当指出，正是在这一理论中，交换子定理的两个特殊情形第一次出现。第一个出现在 I. Schur 的学位论文 {[}44{]} 中，该论文（正是通过在一个张量空间的自同态环中的交换）把线性群的表示与对称群的表示联系起来。第二个出现在 Schur 1905 年的工作 {[}46{]} 中，他在其中证明，在域 C 上与一个不可约表示的所有矩阵都交换的矩阵必是 \(I\) 的标量倍（这个结果也可以由伯恩赛德定理推出）。

剩下的工作是把这些理论的共同基础明确地辨认出来：这是 1925--1933 年间以 E. Noether 和 E. Artin 为核心的德国学派所完成的工作，现代代数正是在这一时期创立的。早在 1903 年，H. Poincaré 就在一篇关于线性微分方程的代数积分的论文 {[}42{]} 中，在一个代数里定义了左理想、右理想以及极小理想的概念。他还观察到，在半单代数中，每个左理想都是它与各个单分量的交的直和，并且在 n 阶矩阵代数中，极小理想的维数是 n；但他的工作没有引起代数学家的注意。\(^{(11)}\) 1907 年，Wedderburn 重新定义了代数的左理想与右理想，并证明了它们的若干性质（特别是证明了根是最大的幂零左理想（{[}51{]}，第 113--114 页））。

但直到 1927 年，这些概念才在代数理论中得到本质性的使用。\(^{(12)}\) 通过对此前散见于各处的证明方法加以推广，\(^{(13)}\) W. Krull 在 1925 年 {[}28{]} 与 E. Noether 在 1926 年 {[}33{]} 引入并系统地使用了极大条件与极小条件。前者用它们把关于有限群分解为不可分解群的直积的雷马克定理推广到带算子的交换群（这一概念就是他在此时定义的）（参见 VIII，第 37 页，定理 2），后者则在刻画戴德金环时使用了这些条件。1927 年，E. Artin {[}1{]} 把同一思想应用于非交换环，说明了如何通过对极小理想的系统研究，把韦德伯恩定理推广到其左理想同时满足极大条件与极小条件的所有环。\(^{(14)}\)

在另一方面，1926 年，Krull {[}29{]} 建立了带算子的交换群的概念与群的线性表示概念之间的联系。这一观点被 E. Noether 在 1929 年一篇奠基性的工作 {[}34{]} 中推广到代数并得到详细发挥。鉴于所引入思想的重要性以及阐述的清晰，这篇著作理应与施泰尼茨关于交换域的论文并列为现代代数的支柱之一。\(^{(15)}\)

最后，在始于 1927 年的一系列工作（{[}36{]}、{[}2{]} 与 {[}35{]}）中，E. Noether 与 R. Brauer（1929 年又有 A. Albert 与 H. Hasse 加入其中）在 Wedderburn 与 Dickson 中断之处重新展开对非交换体的研究。虽然他们的结果中最本质的部分在于对布饶尔群、特别是代数数域上的布饶尔群的深入研究，但正是在这些工作中，交换子定理以及单代数的中和域的概念、它与极大交换子域的关系得到了澄清。最后，1927 年，Skolem 刻画了单环的自同构 {[}49{]}，所得的定理几年后又由 E. Noether {[}35{]} 与 R. Brauer {[}2{]} 重新得出。

就这样，到 1934 年，单环与半单环的初等理论已经或多或少达到了它的最终形式（关于这一理论当时状况的综述，见 {[}16{]}）。
